% Candidate section owned by finite_field_criteria. Root renumbers AC.*.
\subsection{Rational torsion and the relative class number}

We use ``cover'' below to mean a connected, geometrically connected finite
\'etale cover.  A witness in criterion (R) is required to have degree greater
than one.  This requirement is missing from the wording of L1 and A3: with
$G=1$, the numerical condition is automatic, while the proposed contradiction
has no nonidentity deck transformation to use.

\begin{claim}{AC.1}{Proved}
Let $L$ be a finite field of characteristic $p$, let $C/L$ be smooth proper
of genus $g\ge2$, and let $P,Q\in C(L)$ be distinct points with conjugate
sections in the geometrically prime-to-$p$ arithmetic fundamental group.
Let $\pi:D\to C$ be an \'etale Galois cover with constant deck group $G$,
$p\nmid |G|$, and let $P'\in D(L)$ lie over $P$.  There is a unique
$Q'\in D(L)$ above $Q$ with
\[
                 T=[Q'-P']\in J_D(L)[p^\infty].
\]
It is nonzero.  If $|G|$ is odd, every $1\ne\delta\in G$ satisfies
$\delta T\ne T$.  For any homomorphism of abelian varieties
$h:J_D\to B$ over $L$ with $B(L)[p]=0$, we have $h(T)=0$.
\end{claim}
\begin{proof}
We check precisely the part of AP-A.3 used here.  With the rational point
$P$ as base point, $\pi$ corresponds to an open subgroup of the arithmetic
extension containing the chosen representative of $s_P(G_L)$.  Conjugating
the representative of $s_Q$ by a geometric element so that it equals $s_P$
exhibits a fixed element of the finite \'etale fibre over $Q$.  This gives
$Q'\in D(L)$, and the lifted point sections are equal.  The prime-to-$p$
Kummer maps of $J_D(L)$ then kill $[Q'-P']$.  Since $J_D(L)$ is finite,
\[
            \bigcap_{(n,p)=1} nJ_D(L)=J_D(L)[p^\infty].
\]
This proves existence without enlarging $L$.

For uniqueness, we check the proper-curve instance of AP-A.2.  If two such
lifts differ by a deck transformation $a$ of order $m>1$, then
$n(aR-R)$ is principal for some power $n$ of $p$.  Normalize a function
$f$ with this divisor so that its norm for the cyclic action of $a$ is
one.  Hilbert 90 writes $f=a(u)/u$.  Hence
$\operatorname{div}(u)-nR$ is an invariant divisor of degree $-n$.
Every orbit has length $m$, since the action is free on the proper curve.
Consequently $m\mid n$, a contradiction.  The Abel map of a curve of
positive genus is injective on geometric points: a principal divisor
$Q'-P'$ for distinct points would give a degree-one map to $\mathbf P^1$.
Thus $T\ne0$.

We also check AP-A.C2.7 in the full stated characteristic range.  If
$\delta T=T$, then the distinct effective divisors
$\delta Q'+P'$ and $Q'+\delta P'$ are linearly equivalent.  They have no
common term, because their terms above $P$ and $Q$ are disjoint and a
nonidentity deck transformation has no fixed point.  They give a degree-two
map to $\mathbf P^1$, so $D$ is hyperelliptic.  Its unique hyperelliptic
involution is central in its automorphism group.  Choose an odd prime-order
element $a\in G$.  Its induced action on the genus-zero quotient has the
same order and has a geometric fixed point.  The fibre above that point
contains at most two geometric points.  An odd prime-order permutation of
this fibre is trivial, so $a$ fixes a point of $D$, contradicting \'etaleness.
This argument works also in characteristics two and three.  Finally $h(T)$
is an $L$-rational point of $p$-power order, and any nonzero such point has a
nonzero multiple of order $p$.  The hypothesis on $B(L)$ therefore kills it.
\end{proof}

\begin{claim}{AC.2}{Proved}
In the situation of AC.1, assume $|G|>1$ is odd.  Put
$h_L(X)=\#J_X(L)$.  If
\[
                       p\nmid h_L(D)/h_L(C),
\]
then $P$ collides with no distinct point of $C(L)$, and hence with no
point rational over any specified subfield of $L$.  The quotient in this
display is a positive integer.  More generally the same conclusion holds
for any prime-to-$p$ group $G$ containing a nontrivial odd-order element,
provided that element acts freely on $D$ as it does for an \'etale cover.
\end{claim}
\begin{proof}
Let $B=(\ker\pi_*)^0$ be the connected kernel of the norm
$\pi_*:J_D\to J_C$.  Both this morphism and the pullback
$\pi^*:J_C\to J_D$ are homomorphisms over $L$ and satisfy
\[
 \pi_*\pi^*=[d],\qquad \pi^*\pi_*=\sum_{\gamma\in G}\gamma,
 \qquad d=|G|.
\]
Thus $J_C\times B\to J_D$, $(x,b)\mapsto\pi^*x+b$, is an isogeny.
Point counts of abelian varieties over a finite field are isogeny
invariants, so $h_L(D)/h_L(C)=\#B(L)$ is an integer.  On the finite
$p$-primary group $M=J_D(L)[p^\infty]$, division by $d$ is unambiguous.
The endomorphisms
\[
             e=d^{-1}\sum_{\gamma\in G}\gamma,
             \qquad 1-e
\]
are complementary idempotents.  Their images are respectively the invariant
part and the part mapping into $B[p^\infty]$.  For the second assertion,
$\pi_* (1-e)=0$; any possible component of $\ker\pi_*$ has prime-to-$p$
order, since the isogeny above has kernel killed by a power of $d$.
Alternatively $(1-e)M$ is the image on $p$-primary points of
$[d]-\pi^*\pi_*$, whose connected image is contained in $B$.
The numerical hypothesis gives $B(L)[p^\infty]=0$, so $M=eM$.
A hypothetical collision's $T$ would be fixed by every deck transformation,
contradicting AC.1 for any nonidentity odd-order element.

The rigid target here is the new abelian factor $B$ with
$B(L)[p]=0$.  The proof uses rational torsion over this exact $L$;
it makes no assertion that $B$ has geometric $p$-rank zero.
\end{proof}

\begin{claim}{AC.3}{Proved}
For a fixed $C/L$ and $P\in C(L)$, the following are equivalent:
\begin{enumerate}
\item There is a nontrivial constant odd-order prime-to-$p$ Galois
$P$-split \'etale cover satisfying criterion (R).
\item There is such a cover with cyclic group of odd prime order.
\end{enumerate}
The extension field $L$ need not be changed.  Consequently nonabelian
odd-order groups and noncyclic abelian groups cannot enlarge the set of
$(C,P,L)$ reached by the whole-relative-class-number criterion L1.
\end{claim}
\begin{proof}
The odd-order theorem of Feit--Thompson, the unnumbered main theorem in
Chapter I, \S1, printed p.~775, states: ``All finite groups of odd order are
solvable'' \cite{AC_FeitThompson1963}.  We use only its stated hypothesis
that $G$ is finite of odd order.  A nontrivial finite solvable group has a
nontrivial abelianization, hence a quotient $G\twoheadrightarrow
\mathbf Z/\ell$ for an odd prime $\ell\ne p$.  If $N$ is its kernel, let
$E=D/N$.  The intermediate cover $E\to C$ and its deck group are defined
over $L$, and the image of $P'$ is rational, so $P$ splits completely.
For any finite separable map of smooth proper curves over $L$, pullback and
norm give an isogeny decomposition of the source Jacobian with the target
Jacobian as a factor.  Therefore both ratios in
\[
        \frac{h_L(D)}{h_L(C)}=
        \frac{h_L(D)}{h_L(E)}\frac{h_L(E)}{h_L(C)}
\]
are positive integers.  If their product is prime to $p$, so is the second
one.  This proves (1)$\Rightarrow$(2); the reverse implication is immediate.
The argument says nothing about a criterion applied only to a selected
primitive or nonabelian representation factor.
\end{proof}

\subsection{The exact character test and the base-point twist}

\begin{claim}{AC.4}{Proved}
Let $C/L$ be as above, $q=\#L$, $P\in C(L)$, and let
$a:J_C(L)\twoheadrightarrow A$ be a nontrivial finite abelian quotient
of order prime to $p$.  For a nontrivial character
$\psi:A\to\overline{\mathbf Q}^{\times}$ define
\[
 c_d(\psi)=\sum_{E\in\operatorname{Div}^{+}_d(C)(L)}
           \psi\bigl(a([E-dP])\bigr),
 \qquad \kappa=[K_C-(2g-2)P].
\]
Here an effective divisor is a Galois-invariant effective geometric
divisor; it need not be a sum of rational points.  The corresponding
$P$-split Lang cover has twisted polynomial
\[
 L_P(\psi,T)=\sum_{d=0}^{2g-2}c_d(\psi)T^d,
 \qquad c_{2g-2}(\psi)=q^{g-1}\psi(a(\kappa)).
\]
For $g-1\le d\le2g-2$ one has
\[
 c_d(\psi)=q^{d+1-g}\psi(a(\kappa))
                   c_{2g-2-d}(\psi^{-1}).                 \tag{AC.1}
\]
In particular, modulo any prime above $p$,
\[
               L_P(\psi,1)\equiv
                  \sum_{d=0}^{g-1}c_d(\psi).             \tag{AC.2}
\]
For genus two this is exactly
$L_P(\psi,1)\equiv1+\sum_{R\in C(L)}\psi(a([R-P]))$.
\end{claim}
\begin{proof}
The Euler product for the unramified rank-one character, with Frobenius
at a closed point $x$ sent to $a([x-(\deg x)P])$, expands as the displayed
effective-divisor series.  Its coefficient in each degree is $c_d$.
All degree-$d$ line bundles defined over $L$ are parametrized by
$J_C(L)$ after subtracting $dP$.  A line bundle $M$ with
$h^0(M)=r$ has exactly $(q^r-1)/(q-1)$ effective divisors over $L$.
This follows by dividing its nonzero rational sections by $L^\times$;
$\operatorname{Br}(L)=0$ ensures that rational Picard classes are represented
by line bundles.  For $d>2g-2$, Riemann--Roch makes $r=d+1-g$
independent of the class, and summing the nontrivial character over
$J_C(L)$ gives zero.  The series is therefore a polynomial of degree at
most $2g-2$.

For $d\ge g-1$, Riemann--Roch and the change of variables
$M\mapsto\omega_C\otimes M^{-1}$ give
\[
 \frac{q^{h^0(M)}-1}{q-1}
 =q^{d+1-g}\frac{q^{h^0(\omega_C\otimes M^{-1})}-1}{q-1}
   +\frac{q^{d+1-g}-1}{q-1}.
\]
The second term has character sum zero.  The first gives (AC.1), including
its degree-zero coefficient and hence the leading coefficient.  Every
coefficient of degree $d\ge g$ is divisible by $q$, proving (AC.2).
Thus no cohomological interpretation is needed for the genus-two numerical
formula; it is an exact Riemann--Roch calculation.
\end{proof}

\begin{claim}{AC.5}{Proved}
With $A$ as in AC.4 put
\[
 \Theta_{C,P,A}=\sum_{d=0}^{g-1}
   \sum_{E\in\operatorname{Div}^{+}_d(C)(L)}[a([E-dP])]
                      \quad\hbox{in }\mathbf F_p[A].
\]
The full $A$-cover satisfies criterion (R) if and only if
\[
     \bar\psi(\Theta_{C,P,A})\ne0
       \quad\text{for every nontrivial }
          \bar\psi:A\to\overline{\mathbf F}_p^{\times}.    \tag{AC.3}
\]
One can equivalently sum through $2g-2$ in the definition of $\Theta$.
For $A=\mathbf Z/\ell$, with $\ell$ an odd prime different from $p$,
\[
 \frac{h_L(D)}{h_L(C)}
       =\prod_{i=1}^{\ell-1}L_P(\psi^i,1)
       =N_{\mathbf Q(\zeta_\ell)/\mathbf Q}L_P(\psi,1).
\]
If $p$ generates $(\mathbf Z/\ell)^\times$, failure of (R) is equivalent
to all the coefficients of $\Theta$ in $\mathbf F_p[\mathbf Z/\ell]$
being equal.
\end{claim}
\begin{proof}
On a closed fibre of an abelian Galois cover, multiplication by its
Frobenius element has eigenvalues $\psi(a_x)$ in the regular
representation.  Factoring its permutation determinant into these
characters and multiplying over closed points gives
\[
                 Z(D,T)=\prod_{\psi\in\widehat A}L_P(\psi,T).
\]
The trivial factor is $Z(C,T)$.  The common genus-zero denominator of the
two zeta functions cancels, and evaluation of their numerator polynomials
at $1$ gives the relative class-number formula.  By AC.4, reduction of each
factor is its character evaluation on $\Theta$.  Because $p\nmid|A|$,
all characters reduce to distinct characters over
$\overline{\mathbf F}_p$ and exhaust those characters.  The product is
nonzero modulo $p$ exactly when all nontrivial factors are nonzero, which
proves (AC.3).

For prime cyclic $A$, the factors are all Galois conjugates of a fixed
primitive character, giving the norm.  If $p$ is primitive modulo $\ell$,
$1+X+\cdots+X^{\ell-1}$ is irreducible over $\mathbf F_p$.
A polynomial of degree at most $\ell-1$ vanishes at one nontrivial
$\ell$th root exactly when it is a scalar multiple of this polynomial.
The same irreducibility shows that vanishing at one character is equivalent
to vanishing at all nontrivial characters.  This proves the last assertion.
\end{proof}

\begin{claim}{AC.6}{Proved}
For the same fixed character $\psi$ and rational base points $P,P_0$,
\[
 L_{P_0}(\psi,T)
       =L_P\bigl(\psi,\psi(a([P-P_0]))T\bigr).             \tag{AC.4}
\]
The associated covers differ by the constant torsor with arithmetic
Frobenius value $a([P-P_0])$ (up to reversing all Frobenius conventions
consistently).  Thus changing the base point is not innocuous for (R).
For a finite extension $L'/L$, the base change of a $P$-split Lang cover
corresponds to $a\circ N_{L'/L}:J_C(L')\to A$ and remains $P$-split.
If a fixed cover fails (R) over $L$, it fails over every extension $L'$.
If it satisfies (R), it need not do so over every extension.
\end{claim}
\begin{proof}
For every effective divisor of degree $d$,
$[E-dP_0]=[E-dP]+d[P-P_0]$.  Substitution in each coefficient proves
(AC.4).  At closed points this is precisely multiplication of Frobenius
characters by the indicated constant degree character.  Class field
functoriality under constant extension follows directly from the norm of
divisors and the product of Frobenius elements over a closed point; it gives
the asserted norm character.  The cover stays geometrically connected and
has its rational lift, so the resulting quotient is still onto.

For the last assertions use $h(D)/h(C)=\#B(L)$ from AC.2.  A nonzero
$L$-rational $p$-point of $B$ remains rational after every extension.
Conversely a variety with positive geometric $p$-rank acquires nonzero
$p$-torsion over some finite extension even when it has none over $L$.
More precisely, if $F$ is arithmetic Frobenius on $B[p](\overline L)$,
then the criterion survives a degree-$m$ extension exactly when
$\ker(F^m-1)=0$.  This identifies the field dependence rather than
replacing rational torsion by geometric torsion.
\end{proof}

\subsection{The genus-two count and its limits}

\begin{claim}{AC.7}{Proved}
Suppose $g=2$ and $\ell\ne p$ is an odd prime for which $p$ is a primitive
root.  Put $V=J_C(L)/\ell J_C(L)$ and
$N_b=\#\{R\in C(L):[R-P]\bmod\ell J_C(L)=b\}$.
If $P$ collides with a distinct point, then
\[
                      N_b\equiv N_0+1\pmod p
                       \quad(0\ne b\in V).               \tag{AC.5}
\]
If $|V|>\#C(L)$, this implies
$N_0\equiv-1\pmod p$ and $N_b\equiv0\pmod p$ for $b\ne0$.
For one surjection $a:J_C(L)\twoheadrightarrow\mathbf Z/\ell$,
let $H=\ker a$, $N=\#C(L)$.  If
\[
                           N\ge2|H|,
                           \qquad \ell>N,                \tag{AC.6}
\]
then its cover satisfies (R).
\end{claim}
\begin{proof}
Every nontrivial character of $V$ factors through a quotient of order
$\ell$.  Under a collision, AC.2 forces (R) to fail for each such quotient;
AC.5 and the primitive-root hypothesis force every corresponding
nontrivial character evaluation of
$[0]+\sum_R[[R-P]]$ to vanish.  Fourier inversion is valid because
$p\nmid|V|$.  Therefore this group-ring element is a scalar multiple of
$\sum_{b\in V}[b]$, giving (AC.5).  Since $P$ occupies the zero coset,
$|V|>N$ leaves an empty nonzero coset, proving the second assertion.

For the last assertion we verify the geometric input from AP-A.C2.21.
For each nonzero class $h$ in the Jacobian of a genus-two curve, at most
two ordered pairs of distinct geometric points have difference $h$.
Two distinct representations $R-S=R'-S'$ give a degree-two linear
system, hence the unique hyperelliptic pencil, and force
$(R',S')=(\iota S,\iota R)$.  This includes ramification points and
all characteristics.  Thus at most $2(|H|-1)$ rational points share an
$a$-coset with another rational point: each such point occurs as the first
member of an ordered pair with nonzero difference in $H$.
Under (AC.6), at least two points occupy singleton cosets; at most one is
in the zero coset.  There is therefore a nonzero singleton coset.  If (R)
failed, AC.5 applied to this one cyclic quotient, together with an empty
coset (since $\ell>N$), would make each nonzero occupancy divisible by
$p$.  A singleton contradicts this.  Notice that the primitive-root
hypothesis is essential to this argument and is retained explicitly.
\end{proof}

\begin{claim}{AC.8}{Proved}
For any genus $g\ge2$, any finite quotient $A$ of $J_C(L)$, and
$N=\#C(L)$, let $N_b$ count the Abel-image points in the $A$-cosets.
The number of ordered distinct pairs in the same coset satisfies
\[
 \sum_b N_b(N_b-1)
 \le\frac{N^2}{|A|}
       +\frac{|A|-1}{|A|}(2g-2)^2q-N.                     \tag{AC.7}
\]
If the nonzero occupancies are multiples of $p$, the same quantity is
at least $(p-1)(N-N_0)$.  These two estimates alone give a contradiction
only if their displayed numerical bounds conflict; $|A|>N$ by itself
and the inequality $p>(2g-2)^2$ are not the full hypotheses of such a
finite-field contradiction.
\end{claim}
\begin{proof}
Katz, \emph{Twisted L-Functions and Monodromy}, Proposition 9.1.3,
printed p.~185, applies to a smooth proper geometrically connected curve
$C/L$, a point $P\in C(L)$ and any nontrivial character $\rho$ of
$J_C(L)$.  In its $d=1$ case it bounds the absolute value of
$\sum_R\rho([R-P])$ by $(2g-2)\sqrt q$ \cite{AC_Katz2002}.
Its proof checks that the Abel map preserves the geometric abelianized
fundamental group, so the pulled-back nontrivial Lang local system is
geometrically nontrivial; this verifies the needed hypothesis for all
characters factoring through $A$.
Character orthogonality gives exactly
\[
   \sum_bN_b^2=\frac1{|A|}\left(N^2+
             \sum_{\rho\ne1}|\sum_R\rho([R-P])|^2\right).
\]
Apply the bound and subtract $N$ to obtain (AC.7).  For each $b\ne0$
with $N_b>0$ a multiple of $p$, one has
$N_b(N_b-1)\ge(p-1)N_b$.  Sum these inequalities.

In particular the $-N$ in (AC.7) cannot simply be changed to $-q$.
The weaker bound stated in L3 with coefficient $(2g-2)^2-1$ follows
from this argument only if
$N+(2g-2)^2q/|A|\ge q$.  A large-field asymptotic argument also has to
control $N^2/|A|$ and $N_0$ and provide the required quotients over those
fields.  More fundamentally, for $g>2$ the premise on the occupancies
has not followed from a collision: the correct character test AC.2 has
additional effective-divisor terms.  Thus no pair $(p,g)$ is settled
universally by the displayed L3 counting argument without an additional
existence theorem.  AC.7 is the rigorous genus-two sufficient statement.
\end{proof}

\begin{claim}{AC.9}{Refuted: the point-only extension to genus three}
The genus-two expression $1+\sum_{R\in C(L)}\psi([R-P])$ does not test
relative $p$-divisibility in genus three, even for a prime cyclic cover
with $p$ a primitive root.  An explicit example is
\[
        C/\mathbf F_2:\quad y^2+y=x^7+x^3,\qquad P=\infty.
\]
There is a surjection $\chi:J_C(\mathbf F_2)\to\mathbf Z/3$ whose
point-coset counts are $(3,1,1)$ and whose twisted polynomial is
\[
                 L_P(\chi,T)=1+2T+3T^2+4T^3+4T^4.
\]
Its relative class number is $196$, while the point-only expression
reduces to $1$ modulo $2$.  This is a counterexample to exporting that
numerical test; it is not a collision or a counterexample to A1.
\end{claim}
\begin{proof}
The affine curve is smooth because the derivative in $y$ is one.
The pole of the Artin--Schreier function at infinity has order seven,
so the smooth proper normalization has one rational point there and
genus three.  Rational reduced Mumford divisors $(u,v)$ are given by
$u$ monic, $\deg u\le3$, $\deg v<\deg u$ and
$u\mid v^2+v+x^7+x^3$.  Their complete enumeration has $27$ elements.
The class $a=[(0,0)-\infty]$ generates the group of order $27$;
$[(1,0)-\infty]=6a$ and the conjugates have classes $-a$ and $-6a$.
Define $\chi(a)=1\bmod3$.  This specifies the quotient completely.

The five rational points have character exponents $0,1,2,0,0$.
There are no degree-two closed points.  The fifteen effective degree-two
divisors obtained from unordered rational-point pairs have coset counts
$(7,4,4)$.  Consequently, with $\zeta^2+\zeta+1=0$,
\[
 c_1=3+\zeta+\zeta^2=2,
       \qquad c_2=7+4\zeta+4\zeta^2=3.
\]
Since $K_C\sim4\infty$, formula (AC.1) gives $c_3=4$ and $c_4=4$.
Thus $L_P(\chi,1)=14$ and its quadratic norm is $196$.
The point-only expression is $1+c_1=3$, but the correct reduction is
$1+c_1+c_2=6\equiv0$.

The attached program \path{genus3_countercheck.py} enumerates all reduced
divisors, computes the full quotient table, checks closure, inverses,
commutativity and all $27^3=19683$ associativity identities for its exact
Cantor arithmetic, and prints all coefficients.  Its output is
\path{genus3_countercheck_output.json}.  The finite computation is used
for the explicit example, while AC.4 proves the general replacement.
\end{proof}

\begin{claim}{AC.10}{Proved: a finite-prime obstruction for the rational method}
Let $\overline C/\overline{\mathbf F}_p$ have genus at least two and let
$S$ be a finite set of odd primes different from $p$.  Suppose that for
every connected cyclic \'etale cover of degree $\ell\in S$, its new
Jacobian factor has positive geometric $p$-rank.  There is a finite-field
model $C/k$ with $P\in C(k)$ such that no cyclic $P$-split cover of degree
in $S$ satisfies (R), over any finite extension of $k$.
\end{claim}
\begin{proof}
There are finitely many geometric cyclic covers of each fixed degree:
they are determined by the finitely many nonzero homomorphisms from the
finite group $\pi_1(\overline C)^{\mathrm{ab}}/\ell$ to $\mathbf Z/\ell$,
up to changing the deck generator.  Choose a geometric point $P$ and a
lift to each cover.  Descend all these curves, lifts, maps and deck actions
to one finite field $k$.  Enlarge $k$ finitely so that the new factor of
each cover has a nonzero rational point of order $p$; such a point exists
by the rank hypothesis, and there are only finitely many factors.

After any extension $L/k$, a constant cyclic torsor with one of the
specified geometric covers differs from its chosen model by a constant
$\mathbf Z/\ell$-torsor.  Requiring a rational point above $P$ makes this
twist trivial: evaluate the difference character on the section of $P$.
It is therefore the base change of that chosen $P$-split model (up to
changing the deck generator).  The already rational $p$-point persists,
so the relative class number is divisible by $p$.
\end{proof}

\subsection{The primitive rational factor}

\begin{claim}{AC.11}{Proved}
Let $C/L$ be smooth proper of genus $g\ge2$ and not geometrically a
separable double cover of a genus-one curve.  Let $D\to C$ be a
connected \'etale cyclic cover of odd degree $n>1$ prime to $p$, with
constant group $\langle\sigma\rangle$ and $P'\in D(L)$ over
$P\in C(L)$.  Define the homomorphism over $L$
\[
       A=\prod_{\ell\mid n}(1-\sigma^{n/\ell}):J_D\to J_D,
       \qquad P_n=\operatorname{im}(A),
\]
where the product is over distinct prime factors.  If
$p\nmid\#P_n(L)$, then $P$ has no collision.  Equivalently, a collision
forces $p\mid\#P_n(L)$ on every such cover.  This proves L4 with all
field and nontriviality hypotheses explicit.
\end{claim}
\begin{proof}
AC.1 gives a nonzero rational $p$-primary $T=[Q'-P']$ under a collision.
The numerical hypothesis gives $A T=0$.  We check the precise geometric
endgame of AP-A.C2.20 rather than imposing zero geometric $p$-rank.
Put $r=\omega(n)$ and $d=2^r$.  Expanding $A$ gives a principal divisor
with $d$ distinct simple zeros and $d$ distinct simple poles, all supported
on the two free disjoint deck orbits of $P'$ and $Q'$.  The corresponding
function $f$ has separable degree $d$.  The subset products of the elements
$\sigma^{n/\ell}$ are distinct because their orders are distinct primes.
One has $n>(d-1)^2$: for $r=1,2$ this follows from $n\ge3,15$; for
$r\ge3$ use $n\ge3\cdot5^{r-1}>4^r$.

Over $\overline L$ let $M=\overline L(\gamma f:\gamma\in G)$ and
$e=[\overline L(D):M]$.  Poonen's Proposition 2.1, Lemma 2.2 and
Corollary 2.3 in \emph{Gonality of modular curves in characteristic $p$},
author p.~3, apply over the perfect field $\overline L$ and give
\cite{AC_Poonen2007}
\[
                     e\mid d,\qquad g(M)\le(d/e-1)^2.
\]
For completeness, the induction behind the assertion is this.  If an
intermediate field has index $d/s$ and genus at most $(s-1)^2$, adjoining
one more degree-$d$ function with relative degree $a$ gives, by
Castelnuovo--Severi, genus at most
\[
        a(s-1)^2+(a-1)(as-1)
          =(as-1)^2-as(a-1)(s-1)\le(as-1)^2.
\]
This also handles a compositum with more than two generators.
The alternative $e=1$ contradicts
$g(D)=n(g-1)+1\ge n+1>(d-1)^2$.  Thus $e>1$ is a power of two.
Let $Z$ be the smooth curve with field $M$.  Its induced $G$-action is
faithful because the kernel order divides both $e$ and $n$.  It is free:
a prime-order stabilizer of odd order would partition the fibre divisor
of $D\to Z$ into orbits of that prime length, forcing that prime to divide
$e$.  Riemann--Hurwitz gives
$g(Z)-1=n(g(Z/G)-1)$, and the genus bound forces $g(Z)=1$.

Write $f=h\circ v$, $v:D\to Z$ of degree $e$.  Over a zero of $h$, the
fibre of $v$ is reduced, since the zero divisor of $f$ is reduced.  It
lies in the two original deck orbits, each of which maps injectively to
$Z$ by the free action there.  Hence $e\le2$, and so $e=2$.
The separable quadratic involution of $D\to Z$ commutes with $G$, because
$G$ preserves its fixed field and a group of order two has no nontrivial
automorphism.  Taking the quotient by $G$ gives a separable degree-two
map $C\to Z/G$, whose target has genus one.  This contradicts the
hypothesis on $C$.

Every vanishing used before this last geometric contradiction holds over
the original finite field $L$.  The target is the rationally
$p$-torsion-free primitive factor $P_n$, and the remaining argument is
geometric.  Thus the ``controlled-field endgame'' set aside in R8 is not
an additional gap in this conditional rational criterion; the outstanding
issue is constructing a suitable pointed cover over a field on which the
required rational torsion vanishes.
\end{proof}

\begin{claim}{AC.12}{Open}
The nontrivial-cover existence assertion A3 has not been proved or refuted
by AC.1--AC.11.  In its precise reduced form, it asks whether for every
$C/k$ of genus at least two over a finite field and every $P\in C(k)$
there are a finite extension $L/k$, an odd prime $\ell\ne p$ and a
surjection $J_C(L)\to\mathbf Z/\ell$ such that all its nontrivial
twisted values $L_P(\chi^i,1)$ are $p$-adic units.  This is equivalent to
the version allowing every nontrivial odd-order Galois group by AC.3.
\end{claim}
\begin{proof}[Record of the unresolved step]
AC.4--AC.5 provide an exact test and AC.7 supplies a genus-two sufficient
numerical condition.  They do not construct primes with the required
quotient size and primitive-root property, or a substitute for those
conditions.  In higher genus, replacing the effective-divisor test by
point counts is invalid by AC.9.  The random-character estimate in L2 is
not a theorem: character values for a fixed curve are constrained,
and no independence or equidistribution statement with those parameters
has been established.  AC.3 rules out the suggested enlargement to
nonabelian odd-order groups for this particular criterion.  AC.11 gives
a weaker primitive-factor target but does not construct it.  These are
failures of the attempted methods, not counterexamples to the existence
assertion or to A1.
\end{proof}

\begin{claim}{AC.13}{Proved}
Fix any prime $p$, any genus $g\ge2$, and any finite set $S$ of odd
primes different from $p$.  There are a finite field $k$ of characteristic
$p$, a smooth proper geometrically connected ordinary genus-$g$ curve
$C/k$, and $P\in C(k)$, such that no $P$-split cyclic cover of degree
in $S$ satisfies (R), over any finite extension $L/k$.  By AC.3 the same
holds for all odd-order Galois covers whose group has prime divisors only
in $S$.  In particular there is no bound depending only on $(p,g)$ for
the least prime of a rational L1 witness, even if every $(C,P)$ eventually
has some witness.
\end{claim}
\begin{proof}
Ozman--Pries, Theorem 1.1, printed p.~456, assumes $\ell\ne p$ prime,
$g\ge2$ and $0\le f\le g$, with $f\ne0$ when $g=2$.
For every irreducible component $S_f$ of the $p$-rank-$f$ stratum,
its inverse image in the finite cover-moduli space is irreducible and
has generically ordinary Prym \cite{AC_OzmanPries2019}.  Take $f=g$;
all hypotheses then hold.  Ordinariness is an open condition, and the
forgetful map for covers is finite, as stated in their \S2.2.
The image of the closed nonordinary-Prym locus is therefore closed and
omits the generic point.  Thus on a dense open of the ordinary
curve-moduli component, every degree-$\ell$ cyclic cover has ordinary
Prym, of positive dimension $(\ell-1)(g-1)$.

Intersect these dense opens for the finite set $S$.  The intersection
is nonempty.  An \'etale scheme atlas has a closed point over a finite
extension of $\mathbf F_p$ in that open; after a further finite extension
it supplies an actual curve and a rational marked point.  This gives the
geometric hypothesis of AC.10, which then gives the desired finite-field
model and excludes those primes over every extension.  An odd-order
Galois witness supported in $S$ would have a prime-degree witness with
prime in $S$ by AC.3, and is therefore excluded too.
\end{proof}

\subsection{A semilinear obstruction and the least prime for the two hard classes}

\begin{claim}{AC.14}{Proved}
Let $k$ be a finite field of characteristic two, $C/k$ a genus-two curve,
and $P\in C(k)$.  Let $\ell$ be an odd prime with
$\operatorname{ord}_{\ell}(2)=\ell-1$.  Suppose that every geometric
cyclic degree-$\ell$ \'etale cover has ordinary new Jacobian factor.
Suppose further that, for some $d>0$, the action of arithmetic Frobenius
$F^d$ on $J_C[\ell](\overline k)$ is scalar $2^j$, with
$\gcd(j,\ell-1)=1$, and that every eigenvalue of $F$ on this space has
exact order $d(\ell-1)$.  Then no constant-deck cyclic degree-$\ell$
$P$-split cover satisfies (R) over any finite extension of $k$.
\end{claim}
\begin{proof}
Fix a geometric cyclic cover, equivalently a quotient of the geometric
abelianized fundamental group modulo $\ell$.  Its kernel is stable under
$F^d$, because that action is scalar.  Using the section of the actual
point $P$, the open subgroup generated by this kernel and $s_P(G_{k_d})$
constructs a pointed model $D_0\to C_{k_d}$ over the degree-$d$ extension
$k_d/k$.  Its deck group need not yet be constant; Frobenius acts on it
by the scalar $2^j$.  This construction is legitimate before the deck
group is rational.  The $P$-point fixes the twist.

Let $V$ be the new summand of $J_{D_0}[2](\overline k)$, cut out by the
odd-order averaging idempotent.  Ordinariness gives
$\dim_{\mathbf F_2}V=\ell-1$.  Each nontrivial geometric deck character
occurs once, by the decomposition of
$H^1(D_0,\mathcal O_{D_0})$ into degree-zero character lines and their
one-dimensional first cohomology spaces; the Artin--Schreier exact
sequence transfers this to the ordinary stable part.  Since two is
primitive modulo $\ell$, $V$ is the single nontrivial irreducible
$\mathbf F_2[\mathbf Z/\ell]$-module.  Identify it with
$E=\mathbf F_{2^{\ell-1}}$, on which a deck generator acts by multiplication
by a primitive $\ell$th root $\zeta$.

Let $\Phi$ denote arithmetic Frobenius over $k_d$ on $V$.  The relation
$\Phi\zeta\Phi^{-1}=\zeta^{2^j}$ makes it a semilinear map
\[
                            x\longmapsto a x^{2^j}
                  \qquad(a\in E^\times),
\]
possibly with $j$ replaced by $-j$ under the opposite action convention.
Either exponent generates $\operatorname{Gal}(E/\mathbf F_2)$.
Therefore
\[
                 \Phi^{\ell-1}(x)
                   =N_{E/\mathbf F_2}(a)x=x.
\]
All new two-torsion is rational over $k_{d(\ell-1)}$.

A constant cyclic quotient over $k_m$ requires a nonzero quotient of
$J_C[\ell]$ on which $F^m$ is identity.  The exact eigenvalue-order
hypothesis forces $d(\ell-1)\mid m$.  Its $P$-pointed form is the base
change of the constructed pointed model: two abelian torsor forms with
the same geometric kernel differ by a constant character, and evaluation
on $s_P$ makes that difference zero.  It consequently has nonzero rational
new two-torsion over $k_m$.  AC.2 identifies this with even relative class
number, proving the assertion for all finite extensions, not just the
first field of constant deck transformations.
\end{proof}

\begin{claim}{AC.15}{Proved; independent small-prime check for AP-A.C2.24}
Let $\alpha^2+\alpha+1=0$, $\beta=\alpha^2$, and
\[
 C_0/\mathbf F_4:\quad v^2+v=z+\frac{\beta(z+1)}{z^2+z+\alpha}.
\]
Its geometric Raynaud theta divisor has no nonzero point of order three
or five.  Every connected geometric cyclic \'etale cover of either
prime degree has ordinary Prym.  The same statements hold for the
coefficientwise-square conjugate curve.
\end{claim}
\begin{proof}
Put $h=z^2+z+\alpha$ and $Y=hv$, so that
\[
             Y^2+hY=f=h\bigl(zh+\beta(z+1)\bigr).
\]
There are three simple Artin--Schreier poles, namely infinity and the
roots of $h$.  Let their points on $C_0$ be $W_\infty,W_1,W_2$.
The different formula gives
\[
     \operatorname{div}(dz)=2W_1+2W_2-2W_\infty,
     \qquad B=W_1+W_2-W_\infty.
\]
The line of locally exact differentials is $\mathcal O(B)$, since its
local transition functions square to those of differentials and $dz$
is exact.  Thus theta consists of the classes $[R-B]$.
The curve is ordinary: equivalently its three simple Artin--Schreier
poles give two-rank two, or the two independent logarithmic
Cartier-fixed differentials at these poles do.  The origin is not on
theta.  At any of the three ramification points $R$, the class $[R-B]$
is a nonzero two-torsion class; hence an odd-order theta point must have
$R$ away from these points.

If $\ell[R-B]=0$ for odd $\ell$, a function of divisor
$\ell R-\ell W_1-\ell W_2+\ell W_\infty$, multiplied by
$h^{(\ell+1)/2}$, is a function $F$ satisfying
\[
            \operatorname{div}(F)
                 =\ell R+W_1+W_2-(\ell+2)W_\infty.        \tag{AC.8}
\]
The basis of functions regular away from infinity gives
$F=h c(z)+b(z)Y$.  The factor $h$ in the first coefficient is forced by
vanishing at $W_1,W_2$, where $Y=0$.
Taking the hyperelliptic norm and writing $r=z(R)$ gives, after scaling,
\[
        h(c^2+cb)+b^2\bigl(zh+\beta(z+1)\bigr)
                          =(z+r)^\ell.                   \tag{AC.9}
\]
The odd pole order in (AC.8) determines the degree and leading coefficient
of $b$, so none of the following normalizations drops a possible solution.

For $\ell=3$ we have $c=a$ constant and $b=1$.
Comparison of the four coefficients in (AC.9), with $s=a^2+a$, gives
\[
      r=1+s,\qquad r^2=r,\qquad r^3=\alpha r+1.
\]
If $r=0$ the last equation reads $0=1$; if $r=1$ it reads
$1=\alpha+1$.  Both are impossible.

For $\ell=5$ write $c=az+b_0$ and $b=z+d_0$.
Set $A=a^2+a$, $B_0=ad_0+b_0+d_0^2$ and
$C_0=b_0^2+b_0d_0$.  Coefficient comparison gives
\[
 \begin{gathered}
 r=1+A,\qquad B_0=r,\qquad C_0=\beta r+1,\\
 r^4=r+1+\beta d_0^2,\qquad
 r^5=r+\alpha+\beta d_0^2.
 \end{gathered}
\]
Consequently, writing
\[
 \begin{aligned}
 R_5(r)&=r^5+r^4+\beta,\qquad
 S(r)=\alpha(r^4+r+1),\\
 T(r)&=S(r)+r,\\
 U(r)&=r^2+(r+1)S(r)+S(r)^2+\beta r+1,
 \end{aligned}
\]
we must have
\[
                 R_5(r)=0,\qquad d_0^2=S(r),
                 \qquad d_0T(r)=U(r).
\]
Squaring the last equality forces
$U(r)^2+S(r)T(r)^2=0$.  These last two polynomial conditions have
no common root: exact Euclidean division over $\mathbf F_4$ gives
\[
              \gcd\bigl(R_5,U^2+ST^2\bigr)=1.             \tag{AC.10}
\]
Indeed the obstruction reduces modulo $R_5$ to
$R_4=\beta+\beta r+\alpha r^2+r^3+\alpha r^4$, and the exact identity is
\[
 (1+\beta r+\alpha r^2+\beta r^3)R_5
       +(\beta+\alpha r^3+\alpha r^4)R_4=1.
\]
The attached program \path{least_prime_semilinear.py} verifies both the
reduction and this B\'ezout identity, as well as the unreduced identity.

Finally we check the use of R5.  For each nontrivial degree-zero character
line $L^i$ of a cyclic cover, $H^1(C_0,L^i)$ has dimension one.
The locally exact differential sequence identifies the kernel of its
Frobenius arrow to $H^1(C_0,L^{2i})$ with
$H^0(\mathcal O(B)\otimes L^i)$.  The absence of theta points at orders
three and five makes every such arrow bijective, so the whole new
coherent cohomology is Frobenius-stable.  The Artin--Schreier exact
sequence identifies its stable dimension with the new two-rank.
This proves ordinariness of every cover at those orders.  Squaring all
coefficients preserves the proof and gives the conjugate curve.
\end{proof}

\begin{claim}{AC.16}{Proved}
For $C_0$ in AC.15, its conjugate, and every original point
$P\in C_0(\mathbf F_4)$, no cyclic $P$-split cover of degree three or
five satisfies (R) over any finite extension of $\mathbf F_4$.
Thus whenever a degree-seven witness is exhibited for such a point,
seven is its least possible odd prime, with the quantifier over all
finite extensions included.
\end{claim}
\begin{proof}
Direct point counts give $N_4=7$, $N_{16}=15$, hence the characteristic
polynomial of fourth-power Frobenius is
\[
                  P(X)=X^4+2X^3+X^2+8X+16.
\]
The attached certificate checks these counts by direct finite-field
arithmetic and verifies the following exact polynomial identities:
\[
 \begin{array}{c|c|c|c}
 \ell&d&X^d\bmod(P,\ell)&
      \text{order of every root of }P\bmod\ell\\\hline
 3&5&2&10\\
 5&13&2&52.
 \end{array}
\]
For the last column, it checks $P\mid X^n-1$ and
$\gcd(P,X^m-1)=1$ for every proper divisor $m\mid n$, with
$n=10$ or $52$.  The polynomials are separable, since these $n$ are
prime to the corresponding $\ell$.  Cayley--Hamilton therefore supplies
the scalar Frobenius identity on the actual $\ell$-torsion, without
requiring a cyclic integral lattice.  Two is primitive modulo both three
and five.  The rank hypothesis of AC.14 was independently proved in
AC.15.  Apply AC.14 with $(\ell,d)=(3,5)$ and $(5,13)$.
The same Frobenius polynomial and theta argument apply to the conjugate.
No assumption that $P$ is a Weierstrass point was made.
\end{proof}

\subsection{The twisted Frobenius determinant, with coefficients specified}

\begin{claim}{AC.17}{Proved}
Let $D\to C$ be the constant abelian $A$-cover of AC.4 over
$L=\mathbf F_q$.  Let $E\subset\overline{\mathbf F}_p$ be a finite field
containing $L$ and all $|A|$th roots of unity, and let
$\bar\psi:A\to E^\times$ be nontrivial.  Deck transformations act on
cohomology by pullback.  Define
\[
 V_{\bar\psi}=
 \left(H^1(D,\mathcal O_D)\otimes_L E\right)_{\bar\psi},
 \qquad
 \Phi=(f\mapsto f^q)^*\otimes_L 1_E.
\]
Here $(f\mapsto f^q)^*$ means the map on coherent cohomology induced by
the $q$th-power map of the structure sheaf over $L$; it is $L$-linear.
Thus $\Phi$ is its \emph{$E$-linear} extension, not the absolute
$q$-semilinear map on coefficients in $E$.
It preserves $V_{\bar\psi}$, which has dimension $g-1$.  With the local
Artin convention of AC.4 and reduction $\psi\mapsto\bar\psi$,
\[
           \overline{L_P(\psi,T)}
              =\det_E(1-T\Phi\mid V_{\bar\psi}).        \tag{AC.11}
\]
The inverse-character convention replaces $\bar\psi$ on both sides.
In particular this gives the cohomological form of L2 in every genus,
and in genus two
\[
 \Phi\bigm|V_{\bar\psi}
    =-\sum_{R\in C(L)}\bar\psi\bigl(a([R-P])\bigr).
\]
\end{claim}
\begin{proof}
We use the opened theorem of SGA $4\frac12$, \emph{Fonctions $L$ modulo
$\ell^n$ et modulo $p$}, Theorem 2.2(b), with its hypotheses as given
in \S2.1: a separated finite-type scheme over $\mathbf F_q$, a
constructible complex of finite Tor dimension, and a noetherian reduced
coefficient ring of characteristic $p$.  It identifies the Euler product
with the alternating determinant on compactly supported \'etale
cohomology.  Formula (2.4.1) states the consequence for a constructible
sheaf over a coefficient field \cite{AC_Deligne1977}.
Take that field to be $E$, and the sheaf to be the rank-one character
summand of $\pi_*E$.  This meets every hypothesis: $C$ is smooth proper,
the cover is finite \'etale, and the sheaf is locally constant of rank
one.  In the fibre over a closed point $x$, arithmetic Frobenius
translates the torsor by $a([x-(\deg x)P])$.  Its inverse acts on the
space of functions with deck-pullback character $\bar\psi$ by the scalar
$\bar\psi(a([x-(\deg x)P]))$.  Its Euler product is therefore exactly
the reduction of the polynomial in AC.4, with no unnoticed character
inversion.

The theorem makes this reduction the determinant of geometric Frobenius
on the $\bar\psi$-summand of $H^1_{\mathrm{\acute et}}
(D_{\overline L},\mathbf F_p)\otimes_{\mathbf F_p}E$.
There is no degree-zero factor because $D$ is geometrically connected
and $\bar\psi$ is nontrivial.  There is no higher factor: the
Artin--Schreier sequence on the proper curve, together with the
surjectivity of $F-1$ on a finite-dimensional semilinear vector space
over an algebraically closed field, gives
$H^i_{\mathrm{\acute et}}(D_{\overline L},\mathbf F_p)=0$ for $i\ge2$.

Katz, SGA 7 II, Expos\'e XXII, Proposition 2.2.5, gives the functorial
isomorphism between the stable part of coherent cohomology (after scalar
extension) and characteristic-$p$ \'etale cohomology.  It identifies the
linear extension of coherent $q$th-power Frobenius with the inverse of
arithmetic Galois Frobenius; Proposition 2.2.4 records equality of their
characteristic polynomials \cite{AC_Katz1973}.  Because the isomorphism
is functorial for deck automorphisms, apply the idempotent
\[
          e_{\bar\psi}=|A|^{-1}
              \sum_{\gamma\in A}\bar\psi(\gamma)^{-1}\gamma^*.
\]
It gives the same comparison on this character summand.  The nilpotent
coherent part contributes the determinant $1$, so the stable part can
be replaced by all of $V_{\bar\psi}$.  This proves (AC.11), including its
Frobenius convention.

Finally the character summand of $\pi_*\mathcal O_{D_E}$ is a nontrivial
torsion line bundle of degree zero on $C_E$.  Its zeroth cohomology
vanishes, and Riemann--Roch gives first-cohomology dimension $g-1$.
For $g=2$ compare the coefficient of $T$ in (AC.11) with AC.4; it is
minus the scalar $\Phi$, which proves the last formula.
\end{proof}

% AC.18--19 deliberately record a broader criterion than the constant-deck
% convention of the prompt.  They are placed here only for integration.

\begin{claim}{AC.18}{Proved; broader pointed criterion}
Let $L$ be finite of characteristic $p$, let $C/L$ have genus at least
two, and let $P\in C(L)$.  Suppose $\pi:D\to C$ is a finite \'etale
cover, geometrically connected and geometrically Galois, of nontrivial
odd degree prime to $p$.  Assume only that a point $P'\in D(L)$ above
$P$ is specified; the geometric deck group need not be constant over
$L$.  If $p\nmid h_L(D)/h_L(C)$, then $P$ has no collision.

This condition genuinely enlarges the constant-deck criterion over a
fixed field.  The reduction to a prime cyclic quotient in AC.3 is not
asserted for this enlargement: a prime cyclic quotient of the geometric
group need not be stable under arithmetic Frobenius.
\end{claim}
\begin{proof}
A hypothetical collision still gives the rational, nonzero, $p$-primary
class $T=[Q'-P']$ over this very $L$.  The lift argument of AC.1 uses an
open subgroup containing $s_P(G_L)$ and does not require the arithmetic
cover to be Galois.  Its geometric Galois closure has prime-to-$p$ degree,
as required in R1.  The geometric stabilizer argument in AC.1 does not
require deck transformations to be rational either.

Put $d=\deg\pi$ and $B=(\ker\pi_*)^0$.  The isogeny
$J_C\times B\to J_D$ and the averaging projector
\[
                 e=d^{-1}\pi^*\pi_*
\]
on rational $p$-primary points are defined over $L$.
Over $\overline L$, $\pi^*\pi_*=\sum_{\gamma\in G}\gamma$; arithmetic
Frobenius merely permutes the summands.  Thus the same point-count and
projector proof as in AC.2 gives $(1-e)T=0$, so every geometric deck
transformation fixes $T$.  This contradicts the geometric odd-order
stabilizer statement.  Extending $L$ to make the deck group constant
would be an invalid shortcut: it can introduce rational $p$-torsion
and destroy the numerical hypothesis.  AC.19 proves the claimed strict
same-field improvement by a fully specified example.
\end{proof}

\begin{claim}{AC.19}{Proved; explicit nonconstant-deck witness}
Use $\mathbf F_4=\{0,1,\alpha,\alpha+1\}$ with
$\alpha^2+\alpha+1=0$, and let $C_0$ be the curve of AC.15.
Its hyperelliptic quadratic twist
\[
 C_0^{\mathrm t}:\quad
 v^2+v=z+\frac{(\alpha+1)(z+1)}{z^2+z+\alpha}+\alpha
\]
has $\#J_{C_0^{\mathrm t}}(\mathbf F_4)=8$.
At its point $P=\infty$ there is a pointed geometrically cyclic
degree-seven \'etale cover over $\mathbf F_4$ with nonconstant deck
group and relative class number $8051=97\cdot83$, hence odd.
No nontrivial constant odd-order $P$-split Galois cover over
$\mathbf F_4$ satisfies the hypotheses of L1, irrespective of its
relative class number.  This last assertion concerns the fixed field;
it makes no exclusion over its extensions.
\end{claim}
\begin{proof}
Write $h=z^2+z+\alpha$, $Y=hv$.  For $C_0$ and $C_0^{\mathrm t}$ the
polynomials on the right of $Y^2+hY=f$ have ascending coefficient lists
\[
 (1,1,0,\alpha,0,1),\qquad
 (0,1,\alpha,\alpha,\alpha,1),
\]
since twisting adds $\alpha h^2$.  Direct counts give $(N_4,N_{16})$
equal to $(7,15)$ and $(3,15)$ respectively.  Thus the twisted
Frobenius characteristic polynomial is
$X^4-2X^3+X^2-8X+16$, whose value at one is eight.
These counts and both Jacobian orders are independently enumerated by
\path{nonconstant_deck_twist.py}.

Define the original degree-seven cover $D\to C_0$ as the quotient of
the Lang cover, based at $P=\infty$, by the unique order-four subgroup
$H=7J_{C_0}(\mathbf F_4)$.  Let $P'=(P,0)$ be its distinguished point.
The hyperelliptic involution $\iota$ of $C_0$ satisfies
$[\iota R-P]=-[R-P]$.  Inversion on the Jacobian commutes with the Lang
map and preserves $H$.  Consequently it gives a lift
\[
                 \widetilde\iota:(R,z)\longmapsto(\iota R,-z)
\]
to $D$, defined over $\mathbf F_4$, of order two, fixing $P'$ and
conjugating every deck transformation to its inverse.
Twist the pair $(D,C_0)$ by this pair of involutions and the nontrivial
quadratic character of $G_{\mathbf F_4}$.  It gives
$D^{\mathrm t}\to C_0^{\mathrm t}$; because $P'$ is fixed by the lift,
it remains rational.  Arithmetic Frobenius acts by inversion on the
geometric deck group $\mathbf Z/7$, so that group is nonconstant.

For completeness we compute the new numerator rather than infer its
parity from an isomorphism after extension.  Direct enumeration of the
original quotient gives rational-point occupancies
$(1,1,0,2,2,0,1)$, with the point at infinity in the zero coset.
For a primitive seventh root $\zeta$, put $t=\zeta+\zeta^{-1}$.
Then
\[
 t^3+t^2-2t-1=0,
 \qquad S=1+\zeta+\zeta^{-1}+2(\zeta^3+\zeta^{-3})
           =3-t-2t^2,
 \qquad L_{\chi}(T)=1+ST+4T^2.
\]
The involution interchanges the $\chi$ and $\chi^{-1}$ cohomology
summands and commutes with the original Frobenius.  On their direct
sum, twisted Frobenius therefore has determinant
\[
       \det(1-T\widetilde\iota F)
              =L_\chi(T)L_\chi(-T).
\]
Taking the product over the three inverse pairs, or the norm from the
real cubic cyclotomic field, gives the relative numerator
\[
 1+10T^2+65T^4+279T^6+1040T^8+2560T^{10}+4096T^{12}.
\]
At one it is $8051$.  Equivalently, the norms of $5+S$ and $5-S$ are
$97$ and $83$.  The attached exact program verifies the occupancy
vector, the cubic identity, the complete relative polynomial and its
value.  Thus AC.18 applies over $\mathbf F_4$ itself.

Finally any nontrivial constant odd-order pointed Galois cover would,
by the group-theoretic step of AC.3, yield a constant prime cyclic
$P$-split cover.  Class field theory would give a nonzero odd-prime
quotient of $J_{C_0^{\mathrm t}}(\mathbf F_4)$, which has order eight.
This is impossible, proving the last assertion.
\end{proof}

\begin{claim}{AC.20}{Proved; simple Frobenius quotient and exact descent power}
Let $D\to C$ be a cover satisfying the broadened hypotheses of AC.18
over a finite field $L$, and suppose $p\nmid h_L(D)/h_L(C)$.
There is an intermediate pointed geometrically Galois cover $E\to C$
over the same $L$, with geometric deck group
\[
                    V\simeq(\mathbf Z/\ell)^r,
             \qquad \ell\text{ odd},\quad \ell\ne p,
             \quad 1\le r\le2g,
\]
such that arithmetic Frobenius acts irreducibly on $V$ and
$p\nmid h_L(E)/h_L(C)$.  Let $m$ be the order of this action, so
$m\mid\ell^r-1$.  Write $L_m/L$ for the degree-$m$ extension.  Then
the deck group becomes constant over $L_m$ and
\[
             \frac{h_{L_m}(E)}{h_{L_m}(C)}
                  =\left(\frac{h_L(E)}{h_L(C)}\right)^m. \tag{AC.12}
\]
Consequently there is a constant cyclic degree-$\ell$ pointed witness
for L1 over $L_m$.

For a fixed pair $(C/k,P)$, existence of a witness after some finite
extension is therefore equivalent for the broadened criterion AC.18
and for the original constant-deck criterion A3.  Nonconstant deck
groups can lower the field of a witness, as in AC.19, but do not enlarge
the class of pairs reached when all finite extensions are allowed.
\end{claim}
\begin{proof}
Let $G$ be the geometric deck group.  The rational point $P'$ gives an
actual Frobenius action on $G$ by comparing the distinguished point in
the geometric fibre; equivalently the associated finite arithmetic
quotient is the semidirect product determined by this section.
Feit--Thompson makes $G$ solvable, so its abelianization $A$ is
nontrivial.  Frobenius acts on $A$.  Choose an odd prime $\ell$ dividing
its order and pass to the nonzero quotient of its $\ell$-primary part
modulo $\ell$.  Choose a maximal proper Frobenius-stable subspace.
The resulting quotient $V$ is a nonzero irreducible
$\mathbf F_\ell[F,F^{-1}]$-module.  Its kernel in $G$ is normal and
Frobenius-stable.  Thus the corresponding geometric intermediate cover
descends to an $L$-cover $E\to C$, and the image of $P'$ is still
rational.  Since $V$ is a quotient of the geometric abelianized
fundamental group modulo $\ell$, its dimension is at most $2g$.

The same Jacobian isogeny argument as in AC.3 makes
$h_L(D)/h_L(E)$ and $h_L(E)/h_L(C)$ positive integers.  Their product is
prime to $p$, so the second is prime to $p$.  An irreducible finite
$\mathbf F_\ell[F,F^{-1}]$-module is the additive group of
$\mathbf F_{\ell^r}$, with $F$ acting by multiplication by some
$a\in\mathbf F_{\ell^r}^{\times}$.  Its order $m$ divides
$\ell^r-1$.  Every nonzero vector, and every nontrivial character in
the dual module, has orbit of length exactly $m$: if $a^d\ne1$,
multiplication by $a^d-1$ is invertible.  The same conclusion holds for
the inverse or dual action, so no Frobenius convention affects it.

To prove (AC.12), choose an auxiliary prime $\nu\ne p$ and a finite
extension $K/\mathbf Q_\nu$ containing the $\ell$th roots of unity.
On the complement of the invariant summand in
$H^1_{\mathrm{\acute et}}(E_{\overline L},K)$, geometric Frobenius
$\mathcal F$ permutes the nontrivial $V$-character summands.  Group them
by Frobenius orbits.  Each orbit has length $m$.  If $W_0$ is one of
its summands and $A_0=\mathcal F^m|W_0$, transport a basis of $W_0$
around the orbit by $\mathcal F$.  In these bases $\mathcal F$ is a
cyclic block matrix with final block $A_0$.  Elementary elimination
gives
\[
 \det(1-T\mathcal F\mid\text{orbit})
       =\det(1-T^m A_0\mid W_0),\qquad
 \det(1-\mathcal F^m\mid\text{orbit})
       =\det(1-A_0\mid W_0)^m.
\]
Multiply over the orbits and evaluate the first identity at $T=1$.
The invariant part is exactly the cohomology of $C$, because averaging
by $|V|$ splits the finite \'etale pullback and trace.  The determinant
of its complement at $1$ is the relative class number, both over $L$
and over $L_m$.  This proves the exact equality (AC.12), without any
assumption that $m$ is prime to $p$ and without estimating eigenvalues.

Over $L_m$ the geometric deck group $V$ is constant and the chosen point
remains rational.  Equality (AC.12) preserves the prime-to-$p$ condition.
Any quotient $V\twoheadrightarrow\mathbf Z/\ell$ gives a constant
prime cyclic intermediate cover; integer relative class-number ratios
again preserve the condition.  This gives the asserted original L1
witness.  The reverse implication between the two existence assertions
is immediate.  Thus unrestricted enlargement of fields would indeed
have been unsafe in AC.18, but this particular extension after the
irreducible-module reduction is justified by (AC.12).
\end{proof}

\begin{claim}{AC.21}{Proved; exact least-prime-five witness}
For the genus-two curve
\[
                  C/\mathbf F_4:\quad y^2+y=x^5+x
\]
and its hyperelliptic quadratic twist, the least odd prime degree of
a constant pointed L1 witness at $P=\infty$, with all finite extensions
allowed, is five.  A witness is defined over $L=\mathbf F_{256}$.
For the specified witness its four nontrivial twisted polynomials are
all
\[
               L_P(\chi^i,T)=1-32T+256T^2=(1-16T)^2,
                 \qquad 1\le i\le4,
\]
and its relative class number is
$225^4=2\,562\,890\,625=3^8 5^8$, which is odd.
The covering curve is the smooth Fermat quintic.
\end{claim}
\begin{proof}
We use the explicit polynomial basis
$L=\mathbf F_2[t]/(t^8+t^4+t^3+t+1)$; an integer denotes the element
whose binary digits are the coefficients in this basis.  Set
\[
            r=18,\qquad a=254,\qquad b=79,\qquad c=44,
            \qquad F=y+a x^2+b x+c.
\]
The exact certificate \path{supersingular_least5.py} verifies
\[
 \begin{gathered}
 r^{16}+r+1=0,
 \qquad b=1+r^4,\qquad a=b^2,\qquad c^2+c=r^5,\\
 F(F+1)=(x+r)^5,
 \qquad R=(r,42)\in C(L),
 \qquad \operatorname{div}(F)=5R-5P.
 \end{gathered}
\]
For the divisor assertion, $F$ has pole order five at infinity because
$y$ has that order and $a x^2+b x+c$ has pole order at most four.
The norm identity shows that its only zero is at the one point $R$
where $F=0$, with multiplicity five; at its conjugate $F=1$.

Define $D\to C_L$ by $z^5=F$.  It is unramified because every valuation
of $F$ is divisible by five.  It is geometrically connected: otherwise
$F$ would be a fifth power and $R-P$ principal, impossible on a positive
genus curve.  The group $\mu_5$ is constant over $L$ since $5\mid255$.
At $P$, the function $u=x^2/y$ is a uniformizer and the residue of
$u^5F$ is one.  Thus the fibre consists of five rational points, and we
choose the one with residue $uz=1$.  This proves the field and pointing
requirements, rather than obtaining them by an uncontrolled extension.
Putting $w=(x+r)/z$ in the norm identity gives
$w^5+z^5=1$, so the smooth projective covering curve is the Fermat
quintic, of genus six.

The same program enumerates all $193$ points of $C(L)$ and evaluates
the Kummer character using $F(Q)^{51}\in\mu_5$.
At $R$ the leading unit is also one, since
$F/(x+r)^5=(F+1)^{-1}$; its character is therefore zero, as is the
character at $P$.  With $12$ as the primitive fifth root in the chosen
field representation, the five occupancies are exactly
\[
                          (13,45,45,45,45).
\]
For each nontrivial fifth-root character the point sum is $13-45=-32$.
The canonical divisor is $2P$, so AC.4 gives precisely the four
polynomials stated in the claim.  Their product at one gives the
relative class number.  The covering curve has $5\cdot13=65$ rational
points, an additional check of the Kummer character evaluation.

We exclude degree three over every extension, not only the fields in
the numerical search.  The locally exact differential line is
$\mathcal O(P)$, because $\operatorname{div}(dx)=2P$.
If its theta divisor contained a nonzero point of order three, there
would be a point $R_0\ne P$ and a function with divisor $3R_0-3P$.
Every function regular off $P$ has the form $A(x)+B(x)y$.
Its two terms have respectively even pole order and pole order
$5+2\deg B$, so a function with pole order at most three lies in
$\langle1,x\rangle$ and cannot have pole order three.
There is therefore no such theta point.  Since
$\kappa_P=[K_C-2P]=0$, AC.22 excludes every cubic witness over every
finite extension.  This proves minimality of five.

For the quadratic twist, choose $d\in\mathbf F_{16}\subset L$ with
$d^2+d=\alpha$, where $\alpha$ is the nontrivial Artin--Schreier
twisting constant in $\mathbf F_4$.  The substitution $y\mapsto y+d$
gives the isomorphism over $L$, preserving $P$.  Pulling back the
specified Kummer function supplies the same pointed cover and the same
polynomials.  The geometric theta-three exclusion and $\kappa_P=0$
also persist, proving minimality for the twist.
\end{proof}

\begin{claim}{AC.22}{Proved; exact all-extension cubic criterion}
Let $C/k$ be any genus-two curve over a finite field of characteristic
two, let $P\in C(k)$, and put
\[
                       \kappa_P=[K_C-2P]\in J_C(k).
\]
Let $\Theta$ be the geometric theta divisor of the line of locally
exact differentials.  A constant cyclic degree-three $P$-split cover
satisfying L1 exists over some finite extension of $k$ if and only if
at least one of the following holds:
\begin{enumerate}
\item $\Theta$ contains a nonzero point of order three;
\item $\kappa_P\notin3J_C(k)$.
\end{enumerate}
In the second case a witness already exists over $k$.
If neither holds, no extension can produce a cubic witness.

More precisely, for an individual cubic cover over $L\supseteq k$
with character $\chi:J_C(L)\twoheadrightarrow\mathbf Z/3$, criterion
L1 holds exactly when its geometric new two-rank is zero or
$\chi(\kappa_P)\ne0$.
\end{claim}
\begin{proof}
Reduce a primitive character to $\psi:\mathbf Z/3\to\mathbf F_4^\times$
and write
\[
                 S=\sum_{R\in C(L)}\psi(\chi([R-P]))
                         \quad\text{in }\mathbf F_4.
\]
The middle-coefficient Riemann--Roch identity of AC.4 is
$c_1(\psi)=\psi(\chi(\kappa_P))c_1(\psi^{-1})$.
Since inversion on $\mu_3$ is squaring and the coefficients of the
point sum are integers, reduction gives the exact equality
\[
                    S=\psi(\chi(\kappa_P))S^2.          \tag{AC.13}
\]
Thus either $S=0$, or $S=\psi(\chi(\kappa_P))^{-1}$.
By AC.4--AC.5 the relative class number is odd precisely when
$1+S\ne0$; the other character has value $(1+S)^2$, so it gives the
same condition.  If $S\ne0$, this condition is exactly
$\chi(\kappa_P)\ne0$.

It remains to identify $S=0$ correctly.  By AC.17, the coherent
$\psi$-eigenspace is one-dimensional and its $q$th-power Frobenius
scalar is $S$.  The other scalar is $S^2$.
If they vanish, the new stable cohomology has dimension zero, hence
the geometric Prym has two-rank zero by Artin--Schreier.  If they do
not vanish, its two-rank is two, so the new factor is ordinary.
This proves the individual-cover assertion.

The locally exact differential sequence used in R5 says that the
Frobenius arrow on $H^1(C,L_0)$ vanishes precisely when
$L_0\in\Theta$, for a nontrivial degree-zero character line.
The two nontrivial character lines of a cubic cover are interchanged
by squaring.  Consequently a geometric cubic cover has new two-rank
zero exactly when its cyclic order-three subgroup meets $\Theta$.
If the first condition of the claim holds, descend that cover and all
its deck transformations to a finite extension and then rationalize
a point over $P$.  Its geometric two-rank zero is preserved by every
extension and twist, so this procedure gives a valid L1 witness.

If the first condition fails, every geometric cubic cover is ordinary.
If $\kappa_P\notin3J_C(k)$, the nonzero image of $\kappa_P$ in the
finite $\mathbf F_3$-vector space $J_C(k)/3J_C(k)$ is detected by a
linear functional to $\mathbf F_3$.  Its Lang cover is a witness by
the individual criterion.  If instead $\kappa_P=3x$ with $x\in J_C(k)$,
the same equality holds in $J_C(L)$ for every finite extension $L$.
Every cubic character over every such $L$ kills $\kappa_P$, and none
can be a witness.  This proves both directions with the extension
quantifier included.
\end{proof}

\begin{claim}{AC.23}{Proved; canonical-class sufficient condition}
Let $C/k$ have genus two in any positive characteristic $p$, and let
$P\in C(k)$.  Suppose $\ell\ne p$ is an odd prime with
$-1\in\langle p\rangle\subset(\mathbf Z/\ell)^\times$.
If $\kappa_P=[K_C-2P]\notin\ell J_C(k)$, there is a degree-$\ell$
constant $P$-split L1 witness over $k$ itself.
\end{claim}
\begin{proof}
Choose a quotient $\chi:J_C(k)\twoheadrightarrow\mathbf Z/\ell$
which does not kill $\kappa_P$.  For any nontrivial character $\psi$
of this quotient, let $S_\psi$ be the point-character sum reduced
modulo $p$.  Choose $j$ with $p^j\equiv-1\pmod\ell$.
The same middle-coefficient symmetry as above gives
\[
          S_\psi=\psi(\chi(\kappa_P))S_{\psi^{-1}}
                =\psi(\chi(\kappa_P))S_\psi^{p^j}.
\]
If the genus-two test failed at this character, $1+S_\psi=0$,
so $S_\psi=-1$ and this equality would force
$\psi(\chi(\kappa_P))=1$.  That contradicts the choice of $\chi$,
because every nontrivial character of a group of prime order is
faithful.  Thus no nontrivial factor vanishes modulo $p$, and AC.5
proves L1.  The assumption is only that inversion is a power of $p$;
it does not require $p$ to generate the whole multiplicative group.
\end{proof}

\begin{claim}{AC.24}{Proved; ramified elliptic-factor obstruction}
Let $L$ be a finite field of characteristic two.  Suppose $C/L$ has
genus at least two and a separable degree-two morphism
$f:C\to E$ to an elliptic curve.  Let $P\in C(L)$ be a ramification
point, and use $f(P)$ as the origin of $E$.
Let $\varphi:E'\to E$ be an \'etale isogeny of odd degree $n>1$
over $L$, with no assumption that its kernel is constant, and let
\[
                         D=C\times_E E'.
\]
Then $D\to C$ is a connected \'etale pointed geometrically Galois
cover with odd geometric deck group, but
\[
                            2\mid h_L(D)/h_L(C).
\]
Consequently a prime cyclic character factoring through this elliptic
quotient cannot be an L1 witness.  The assertion persists over every
finite extension.  Neither L1 nor its nonconstant-deck enlargement
AC.18 can use this elliptic-isogeny pullback.
\end{claim}
\begin{proof}
The cover $E'\to E$ is geometrically connected of odd degree, whereas
$C\to E$ has degree two.  Their function fields are linearly disjoint,
since their intersection degree divides both degrees.  Thus the fibre
product is geometrically connected; it is smooth because $D\to C$
is \'etale.  Put $G=\ker\varphi(\overline L)$, a group of order $n$.
The fibre over $P$ consists of ramification points $R_g$ of the
quadratic map $\rho:D\to E'$, indexed by $B_g=g\in G$.
The point $R_0$ over the origin of $E'$ is rational; the remaining
points are permuted by Galois.  Set
\[
        t=n\left[\sum_{g\ne0}R_g-(n-1)R_0\right]\in J_D(L).
\]
This class is rational because its divisor is Galois-invariant.
Each $n[B_g-B_0]=0$, and the quadratic pullback sends $B_g$ to
$2R_g$, so $2t=0$.  We prove $t\ne0$, including geometrically.
Otherwise there would be a function $u$ on $D_{\overline L}$ with
$\operatorname{div}(u)=n\sum_{g\ne0}R_g-n(n-1)R_0$.
The quadratic involution $\tau$
fixes each ramification point, so $\tau(u)/u$ is a constant $c$.
Applying $\tau$ again gives $c^2=1$, and characteristic two forces
$c=1$.  Thus $u$ is pulled back from a function on $E'_{\overline L}$.
Its valuations at these ramification points must be even, contradicting
the odd valuation $n$ at any $R_g$ with $g\ne0$.
Hence $t$ is nonzero rational
two-torsion.

The deck average of every difference $R_g-R_0$ is zero, so
\[
                          \sum_{\gamma\in G}\gamma(t)=0.
\]
On two-primary torsion the averaging projector is defined because
$n$ is odd.  It kills $t$, putting $t$ in the new summand, and hence in
the connected norm kernel on two-primary torsion as in AC.2.
The projector is defined over $L$ even when its individual summands
are not, as checked in AC.18.
This new abelian factor therefore has an even number of rational
points.  Its point count equals $h_L(D)/h_L(C)$.

Finally, a constant odd prime cyclic unramified cover of $E$ pointed
over its origin is an \'etale elliptic isogeny after using the chosen
point as the origin upstairs.  A character of $J_C(L)$ obtained by
composition with $f_*:J_C(L)\to E(L)$ gives its pullback.  This
establishes the last statement for exactly the covers in question;
it does not exclude characters involving more than one elliptic
factor of $J_C$.
\end{proof}

\begin{claim}{AC.25}{Proved; a fixed-field obstruction to the enlarged criterion}
Let $c\in\mathbf F_4\setminus\mathbf F_2$ and
\[
           C_c/\mathbf F_4:\quad y^2+y=x+\frac{c}{x^2+x},
                   \qquad P=\infty.
\]
There is no pointed geometrically Galois cover of nontrivial odd
degree over $\mathbf F_4$ satisfying the relative-class-number
criterion, even if its geometric deck group is nonconstant or
nonabelian.  Thus the variant of A3 requiring the original field
while allowing nonconstant odd deck groups is false.
This does not refute A3, which allows finite extensions.
\end{claim}
\begin{proof}
Choose $\alpha\in\mathbf F_4$ with $\alpha^2+\alpha=1$ and put
$\beta=\alpha^2$.  The two involutions
\[
  \tau_1(x,y)=(x+1,y+\alpha),\qquad
  \tau_2(x,y)=(x+1,y+\beta)
\]
are defined over $\mathbf F_4$, fix $P$, and have product the
hyperelliptic involution $\iota$.  Their elliptic quotients have
explicit coordinates
\[
 \begin{array}{c|c|c}
 &u&w\\\hline
 E_1&x^2+x&y+\alpha x\\
 E_2&x^2+x&y+\beta x
 \end{array}
 \qquad
 \begin{aligned}
 E_1:&\quad w^2+w=\beta u+c/u,\\
 E_2:&\quad w^2+w=\alpha u+c/u.
 \end{aligned}
\]
Each map has degree two and ramifies at $P$.
Using the images of $P$ as origins, let
$\Phi:E_1\times E_2\to J_C$ be the sum of their Jacobian pullbacks,
and $\Psi:J_C\to E_1\times E_2$ their pair of norms.
The hyperelliptic involution is $[-1]$ on $J_C$, so
$\tau_2=-\tau_1$ there.  Pullback followed by norm is $[2]$;
the cross terms vanish, since each is both invariant and anti-invariant
under $\tau_1$ and a homomorphism of abelian varieties killed by two
must be zero.  Consequently
\[
                   \Psi\Phi=[2],\qquad \Phi\Psi=[2].
\]
These are isogenies whose kernels are killed by two, giving a direct
sum decomposition on every odd Tate module.

Counting the two elliptic curves over $\mathbf F_4$ gives the
unordered pair of point counts $8,4$, hence the Frobenius polynomials
\[
                       X^2+3X+4,\qquad X^2-X+4.
\]
The attached certificate verifies these counts for both values of $c$.
Modulo every odd prime $\ell$, these polynomials are coprime: their
difference is $4X$, while neither has zero as a root.
Every Frobenius-simple quotient of the odd geometric abelianized
fundamental group must therefore be supported in just one of the two
elliptic factors.  A pointed cover with that quotient is the pullback
of the corresponding pointed unramified cover of that elliptic curve.
Such a genus-one cover, with its chosen point used as origin, is an
\'etale elliptic isogeny over $\mathbf F_4$, with no assumption that its
kernel be constant.  The pointing removes a possible constant torsor
twist, so this comparison is over the same field.

If an odd nonabelian or nonconstant witness existed, AC.20 would first
give such a Frobenius-simple elementary abelian quotient witness over
the same field.  AC.24 excludes its elliptic-isogeny pullback, since
both elliptic quotient maps ramify at $P$.  This contradiction proves
the assertion for every odd group, not merely the covers in a finite
search or a multiplication-isogeny family.
\end{proof}

\begin{claim}{AC.26}{Proved; stronger all-extension small-prime and field obstructions}
For the pointed curves $(C_c,P)$ of AC.25:
\begin{enumerate}
\item No cubic or quintic constant L1 witness exists over any finite
extension of $\mathbf F_4$.
\item No nontrivial constant odd-group L1 witness exists over
$\mathbf F_{4^m}$ for any $1\le m\le15$.
\item At $m=16$, the only possible prime of a prime cyclic witness is
$113$. The positive certificate at this degree is supplied in AC.28.
\end{enumerate}
These conclusions about constant covers also hold at every rational
hyperelliptic ramification point of $C_c$.
\end{claim}
\begin{proof}
We first check the geometric small-prime assertions independently of
the previous-return certificates.  Write $h=x^2+x$, $Y=hy$, so
$Y^2+hY=h(xh+c)$.  The locally exact line has divisor
$B=W_0+W_1-P$.  An odd-order theta point is of the form $[R-B]$ with
$R$ away from the three ramification points.  If its order divides
$n=3$ or $5$, multiplication of a principal function by
$h^{(n+1)/2}$ gives a function with divisor
$nR+W_0+W_1-(n+2)P$.

For $n=3$, its pole bound and zeros at $W_0,W_1$ force it, after
scaling, to be $ah+Y$.  Taking its norm and dividing by $h$ gives
\[
             h(a^2+a)+xh+c=(x+r)^3.
\]
Coefficient comparison yields $r^2+r=1$ and $r^3=c$.
The first equation gives $r^3=1$, contradicting $c\ne1$.
Thus theta has no nonzero three-torsion.  Since $K_C\sim2P$,
AC.22 excludes every cubic witness over every extension.

For $n=5$, the same pole and zero conditions give
$h(ax+b)+(x+d)Y$.  Put $s=d^2$ and compare its norm, divided by $h$,
with $(x+r)^5$.  The resulting necessary conditions are
\[
 \begin{gathered}
 Q=r^4+r+1+c=0,\qquad s=r^5/c,\qquad a^2+a=r+1,\\
 E=ad+b=s+r+1,\qquad B_0=b^2+bd=r+1+c,\\
 dE=U,\qquad U=B_0+E^2+(r+1)s.
 \end{gathered}
\]
Hence $U^2+sE^2=0$.
Using $c^2+c=1$, its remainder modulo $Q$ is
$R_2=(c+1)(r^2+r)$.  There is the explicit identity
\[
                    cQ+(c+1)(r^2+r+1)R_2=1.
\]
Both this remainder and the identity are checked over $\mathbf F_4$
in \path{bielliptic_obstructions.py}.  Thus theta has no nonzero
five-torsion either.  By the checked locally exact differential
sequence, every geometric cyclic quintic cover has ordinary new
factor, of two-rank four.

We must still exclude rational quintic witnesses, rather than
substituting geometric ordinariness for that assertion.  From AC.25,
Frobenius on $J_C[5]$ has characteristic polynomial
\[
                   (X-3)^2(X^2-2X-1)\quad\text{over }\mathbf F_5.
\]
On the first primary part it is $3I+N$, with $N^2=0$; on the second
part its cube is $2I$, and every nonzero vector has period twelve.
Apply the same decomposition to a nonzero quotient functional $v$.
Its period $m$ and the action at $d=m/4$ are exhausted by
\[
 \begin{array}{c|c|c|c}
 \text{support of }v&m&d&vF^d/v\\\hline
 \text{first part, }vN=0&4&1&3\\
 \text{first part, }vN\ne0&20&5&3\\
 \text{second part present, }vN=0&12&3&2\\
 \text{second part present, }vN\ne0&60&15&2.
 \end{array}
\]
The computation checks all $624$ nonzero vectors for each of the two
possible Jordan types, in addition to the displayed algebraic proof.
The geometric cover has a pointed model over $\mathbf F_{4^d}$ with
deck Frobenius equal to $2$ or $3$; both generate
$\mathbf F_5^\times$.  On its ordinary new two-torsion, the argument
of AC.14 identifies this Frobenius with a semilinear map of the
one-dimensional $\mathbf F_{16}$-space, whose fourth power is the
norm to $\mathbf F_2$ and therefore the identity.  Thus all new
two-torsion is rational at the first constant-deck field
$\mathbf F_{4^m}$ and remains rational over every later eligible
field.  This excludes every quintic witness.

For the field bound, the isogenies in AC.25 identify the odd-prime
quotients of $J_C(\mathbf F_{4^m})$ with those of the two elliptic
groups.  If a character is supported in just one factor, AC.24
excludes it.  A remaining prime must therefore divide both elliptic
orders.  With traces $t_1=-3,t_2=1$, define
\[
 u_i(0)=2,\quad u_i(1)=t_i,\quad
 u_i(m)=t_i u_i(m-1)-4u_i(m-2),\quad
 N_i(m)=4^m+1-u_i(m).
\]
The exact gcd computation gives, for $1\le m\le15$,
\[
 \gcd(N_1(m),N_2(m))_{\mathrm{odd}}=
 \begin{cases}
 3,&m=4,8,\\
 135,&m=12,\\
 1,&\text{otherwise}.
 \end{cases}
\]
Only three and five occur, and both were excluded over all
extensions.  At $m=16$ the odd gcd is $339=3\cdot113$.
This proves the two field assertions for prime covers and hence,
by AC.3, for all constant odd groups.

Finally any two hyperelliptic ramification points differ by a class
killed by two in the Jacobian.  Every odd-order character over every
extension kills their difference.  The exact twist formula AC.6
therefore identifies their pointed abelian covers and relative
polynomials.  AC.3 reduces the odd-group assertion to these prime
abelian covers, proving the last statement.
\end{proof}


\begin{claim}{AC.27}{Proved; exact Cartier--Mumford rational test}
Let $k=\mathbf F_q$, $q=2^d$, and let
\[
 C/k:\quad Y^2+h(x)Y=f(x),\qquad \deg h\le2,
 \quad f\text{ monic of degree }5
\]
be smooth of genus two, based at its unique point $P=\infty$.
Let $n>1$ be odd and let $a\in J_C(\overline k)$ have exact order
$n$ and satisfy $\sigma_q(a)=[q]a$, where $\sigma_q$ is arithmetic
Frobenius. Choose a finite extension $E/k$ containing $a$ and $\mu_n$.
For $0<i<n$, write the unique reduced Mumford representative of $ia$
as $(u_i,v_i)$, with $u_i$ monic of degree $r_i\in\{1,2\}$, so that
$ia=[D_i-r_iP]$.

For $j=2i\bmod n$, compute $c_i$ as follows.
If $r_i=1$, put $c_i=h_1$, the coefficient of $x$ in $h$.
If $r_i=2$ and $\gcd(u_i,h)$ has degree one, put $c_i=h_2$.
Otherwise $r_i=2$ and $\gcd(u_i,h)=1$; let $v$ be the Cantor
doubling interpolation of degree at most three, with
$u_i^2\mid v^2+hv+f$ and $v\equiv v_i\bmod u_i$.
If the reduced double has degree one, put $c_i=1$.
If it has degree two, let $v_3$ be the cubic coefficient of $v$
and $b_j$ the coefficient of $x$ in $u_j$, and put $c_i=v_3b_j$.
No other case occurs for a nonzero point of odd order. Define
\[
 A_i=\prod_{s=0}^{d-1}c_{2^s i}^{\,2^{d-1-s}}\in E.       \tag{AC.14}
\]
There is a unique constant cyclic $P$-split degree-$n$ cover over $k$
with geometric character subgroup generated by $a$.
Its linear coherent $q$-Frobenius scalar on the summand for $ia$ is
$\sigma_q^{-1}(A_i)$. Consequently its relative class number is odd
if and only if
\[
                         A_i\ne1\qquad(0<i<n).           \tag{AC.15}
\]
For the primitive factor alone use only $\gcd(i,n)=1$.
These formulas refer to the pointing at infinity; another base point
requires the twist of AC.6.
\end{claim}
\begin{proof}
Put $L_i=\mathcal O(D_i-r_iP)$ and rigidify it at $P$ using
$t=x^2/Y$. Monicity of $f$ gives leading residues one for $xt^2$
and $Yt^5$. The differential $dx/h$ has divisor $2P$.
At finite unramified points this follows from $x$ being a local
parameter. At a finite ramification point,
$dx/h=dY/(f'+h'Y)$ is nonvanishing by smoothness.
Thus its canonical divisor is supported at $P$ and has degree two.
A basis of the one-dimensional Serre dual space
$H^0(C_E,\Omega_C\otimes L_i^{-1})$ is
\[
                            \omega_i=u_i\,dx/h.
\]
Indeed $\operatorname{div}(u_i)=D_i+\iota D_i-2r_iP$ and $r_i\le2$.
Let $e_i$ be its dual basis in $H^1(C_E,L_i)$.

Let $g_i$ have divisor $2D_i-D_j-(2r_i-r_j)P$, normalized to
leading residue one at $P$ relative to $t$.
Multiplication by $g_i$ identifies $L_i^2$ with $L_j$ respecting
rigidifications. Absolute Frobenius followed by this identification
is square-semilinear. Cartier is its Serre adjoint:
\[
 \langle g_iF(z),\omega\rangle
             =\langle z,\mathcal C(g_i\omega)\rangle^2.
\]
This identity follows directly in the residue pairing: the coefficient
of $t^{-1}dt$ in $z^2\eta$ is the square of the residue of
$z\mathcal C(\eta)$. Therefore
$\mathcal C(g_i\omega_j)=\sqrt{c_i}\omega_i$ is exactly the assertion
$F(e_i)=c_i e_j$.

We verify all four cases.
If $r_i=1$, the point is unramified, since a ramification point would
give a two-torsion class. Then $D_j=2D_i$, $u_j=u_i^2$, and $g_i=1$.
The identity $\mathcal C(dx/h)=\sqrt{h_1}\,dx/h$ gives $c_i=h_1$.

If $r_i=2$ and $\gcd(u_i,h)=x-w$, write $D_i=W+R$ with $W$
ramified and $R$ unramified. Since $2W\sim2P$, we have $D_j=2R$
and $g_i=x-w$. With $u_i=(x-w)(x-r)$, direct extraction of odd
powers gives
\[
 \mathcal C\bigl((x-w)(x-r)^2dx/h\bigr)
             =\sqrt{h_2}(x-w)(x-r)dx/h.
\]
For this calculation write $(x-w)/h=(x-w)h/h^2$ and use $h(w)=0$:
the square root of the odd part is $\sqrt{h_2}(x-w)$.
Hence $c_i=h_2$. Two ramification points in $D_i$ would give a
two-torsion class and are excluded.

In the remaining cases $\gcd(u_i,h)=1$.
For a degree-two reduced double, Cantor reduction gives
\[
 v^2+hv+f=v_3^2u_i^2u_j,\qquad
                     g_i=(Y+v)/(v_3u_j).
\]
This function has the required divisor and normalized leading residue.
The Artin--Schreier relation for $Y/h$ gives
\[
 \mathcal C((Y+v)dx/h)=v_3\sqrt{b_j}\,u_i dx/h.
\]
Indeed the term $(Y/h)^2dx$ has Cartier zero, and the odd powers of
$f+hv$ are exactly those of $v_3^2u_i^2u_j$, since $v^2$ has only
even powers. Dividing by $v_3$ before Cartier gives
$\sqrt{v_3b_j}\omega_i$.
For a degree-one double, $v$ has degree at most two, the norm
$v^2+hv+f=u_i^2u_j$ is monic of degree five, and
$g_i=(Y+v)/u_j$. The same calculation with $u_j$ monic linear
gives Cartier coefficient one.

Iteration gives $F^d(e_i)=A_i e_{qi}$, with $A_i$ as in (AC.14).
The hypothesis $\sigma_q(a)=qa$ makes coefficientwise conjugation
of the reduced representative $(u_i,v_i)$ equal to $(u_{qi},v_{qi})$.
Consequently $\sigma_q(\omega_i)=\omega_{qi}$ and
$\sigma_q(e_i)=e_{qi}$, with rigidifications unchanged.
The linearized coherent operator is $\sigma_q^{-1}F^d$,
whose scalar is $\sigma_q^{-1}(A_i)$.
This explains the $q$th-root correction; it cannot change whether
the scalar is one.

Finally, $\sigma_q(a)=qa$ is precisely invariance of the associated
geometric homomorphism to the constant group $\mathbf Z/n$.
The factor $q$ records Frobenius on $\mu_n$ in the Kummer
identification. The section of $P$ gives its unique pointed descent.
All multiplication functions were normalized at $P$, so no constant
character twist is omitted. Applying AC.17 to these one-dimensional
character spaces proves (AC.15). Selecting only primitive characters
gives the primitive-factor statement.
\end{proof}


\begin{claim}{AC.28}{Proved; mixed witnesses and exact prime/field minima}
For the pointed split curves $(C_c,\infty)$ of AC.25, the least
odd prime of a constant pointed L1 witness over all finite extensions
is eleven. A degree-eleven witness is defined over $\mathbf F_{4^{60}}$,
and sixty is the least extension degree for that prime.
The least extension degree permitting any constant odd-group witness
is sixteen, where a degree-$113$ witness exists.

For the following eight $\mathbf F_4$ models, write
$c\in\{\alpha,\alpha^2\}$ and $\delta\in\{0,\alpha\}$:
\[
 y^2+y=x+\frac{c}{x^2+x}+\delta,\qquad
 y^2+y=x+\frac{c}{x^2+x+\alpha}+\delta.
\]
Each original rational point has an eleven witness over
$\mathbf F_{4^{60}}$ and a $113$ witness over $\mathbf F_{4^{16}}$.
Whenever the cubic test AC.22 fails, eleven is the least odd prime.
In particular this holds for every rational point of the split
models and for infinity on each quadratic-pole model.
\end{claim}
\begin{proof}
We first specify and independently check the positive certificates.
Use the field
\[
 K=\mathbf F_2[t]/(t^{120}+t^4+t^3+t+1)=\mathbf F_{4^{60}},
\]
with binary polynomial coefficients encoded by integers, and set
\[
 \alpha=931750399618260646238930031330849980.
\]
The model is $Y^2+(x^2+x)Y=x^5+x^3+\alpha x^2+\alpha x$.
The following reduced pair defines the degree-eleven character line:
\[
 \begin{aligned}
 u(x)&=x^2+u_1x+u_0,\qquad v(x)=v_1x+v_0,\\
 u_0&=1286551050156223861353970523064590862,\\
 u_1&=855838071379080258553061447214623248,\\
 v_0&=372422475677710747247258963726964873,\\
 v_1&=880561393008999941135507307093448287.
 \end{aligned}
\]
The certificate verifies irreducibility of the field polynomial,
$\alpha^2+\alpha=1$, validity of the reduced pair, and exact order
eleven of its Jacobian class $a$. Since $q=2^{120}\equiv1\bmod11$,
the descent condition $\sigma_q(a)=qa$ holds over $K$.

All ten reduced multiples and Cartier coefficients are recorded in
\path{bielliptic_mixed_ell11_F4pow60.json}.
The independent verifier \path{verify_mixed_certificates.py} uses
pure Python carryless field arithmetic, with neither the searcher's
compiled field helper nor its Cartier-coefficient routine.
It recomputes every multiple and reconstructs doubling interpolation
from $(f+hv+v^2)/u$ and $h^{-1}\bmod u$.
All ten $q$-Frobenius scalars differ from one. Their product
determinant has the especially compact form
\[
          \overline{P_D(T)/P_C(T)}=1+T^4+T^{10}.
\]
Equivalently, the first scalar has reciprocal minimal polynomial
$1+T^2+T^5$ over $\mathbf F_2$, and its five conjugates each occur
twice. The determinant evaluates to one at $T=1$.
AC.27 therefore proves the odd relative class number over $K$,
with the cover pointed at infinity. This certificate determines
the factors modulo two; their full characteristic-zero coefficients
are not reconstructed here.

For the smaller field degree, the second certificate uses
\[
 E=\mathbf F_2[t]/(t^{224}+t^{21}+t^7+t+1)
\]
and its subfield $L=\mathbf F_{2^{32}}=\mathbf F_{4^{16}}$.
The complete reduced pair and all $112$ multiples are in
\path{bielliptic_mixed_ell113_F4pow16.json}.
They have exact order $113$ and satisfy $\sigma_{2^{32}}(a)=16a$,
where $16\equiv2^{32}\bmod113$.
Thus the constant pointed cover descends to $L$, although its
displayed character line uses the splitting field $E$.
The independent verifier checks all $112$ coefficients, applies
the actual inverse-$q$ conjugation in AC.27, and obtains an
$\mathbf F_2$ determinant with $31$ nonzero coefficients, whose
value at one is one. This proves the claimed witness over $L$.
Together with AC.26 it proves exact minimality of the field degree
sixteen.

AC.26 excludes three and five over every extension, and AC.30
excludes seven. The eleven witness therefore proves exact minimality
of the prime eleven. To minimize its field, the two elliptic factors
have modulo-eleven eigenvalue orders five and sixty respectively:
the first polynomial factors as $(X-3)(X-5)$, and the second has
both roots of exact order sixty. These facts are verified in
\path{bielliptic_obstructions.py}.
An eleven character using only one elliptic factor is excluded by
AC.24; using both requires an extension degree divisible by sixty.
Our witness attains it.

We also check transfer to the stated forms and points.
A quadratic-pole model becomes split after $x=z+s$, with
$s^2+s=\alpha$, followed by an Artin--Schreier constant change of
$y$. Both changes are available over $\mathbf F_{4^4}$ and preserve
infinity. The split quadratic twist is handled by the constant
change alone. This field embeds in both witness fields, of degrees
sixty and sixteen, so the transported witnesses retain exactly
those fields. Coefficientwise squaring transfers the certificates
from $c=\alpha$ to $c=\alpha^2$; no assertion that these two
geometric classes are $\mathbf F_4$-linearly isomorphic is needed.

The original Jacobian orders of these eight forms are $32,12,48,8$,
as given by their two point counts or the root's census. They are
prime to eleven and $113$. Hence any original point difference
$[Q-\infty]$ has order prime to the cover degree, also after transport
to the witness field. Every character of that prime degree kills
the difference. AC.6 proves that the transported cover is split
over every original $Q$, with the same relative polynomials.

For minimality, a hypothetical three, five or seven witness on a
form can be extended to the composite with $\mathbf F_{4^4}$.
The extra field degree is a power of two, which preserves the
numerical condition by AC.29. The isomorphism then transfers it
to the split model. Five and seven are unavailable at all original
points: their differences from infinity have orders dividing
$32,12,48$ or $8$, hence are divisible by those primes in every
relevant rational Jacobian, so their base-point twists are trivial.
For three we use the exact test AC.22. The split untwisted
Jacobians have two-power order, while the three points on each
split quadratic twist are hyperelliptic ramification points;
their cubic test fails. Infinity has $\kappa=0$ on every form.
This proves all minimal-prime assertions in the claim.
\end{proof}

\begin{claim}{AC.29}{Proved; safe extensions of characteristic-power degree}
Let $B/k$ be an abelian variety over a finite field of characteristic
$p$, and let $k'/k$ have degree $p^r$. Then
\[
 B(k)[p]=0\quad\Longleftrightarrow\quad B(k')[p]=0.
\]
In particular a relative-class-number $p$-unit for a fixed finite
separable cover is preserved by every constant extension of
$p$-power degree.
If two pointed curves become isomorphic over a $p$-power extension,
their existence of prime-degree witnesses over arbitrary extensions,
and the set of primes furnishing such witnesses, agree.
\end{claim}
\begin{proof}
On the finite-dimensional vector space
$V=B[p](\overline k)$ let $F$ be arithmetic Frobenius.
Rational $p$-torsion is its fixed space, and
\[
                     F^{p^r}-1=(F-1)^{p^r}.
\]
A power of an endomorphism of a finite-dimensional vector space
is invertible exactly when that endomorphism is invertible.
Thus the two fixed spaces are simultaneously zero.
Apply this to the connected new Jacobian factor to get the
relative-class-number assertion.
For pointed forms, the composite of any witness field with the
given $p$-power isomorphism field has $p$-power degree over the
witness field. The witness survives this extension and transports
across the pointed isomorphism. Apply the same argument in the
other direction.
\end{proof}

\begin{claim}{AC.30}{Proved; global seventh-prime exclusion for the bielliptic classes}
For $(C_c,\infty)$ of AC.25, no constant pointed degree-seven
L1 witness exists over any finite extension of $\mathbf F_4$.
The same holds for its coefficient conjugate and every pointed
form isomorphic to it at infinity over a power-of-two extension.
\end{claim}
\begin{proof}
Modulo seven the Frobenius polynomial from AC.25 is
\[
                  (X-2)^2(X^2-X+4).
\]
On the first primary part write $F=2I+N$, $N^2=0$.
On the second part every nonzero vector has period twenty-four
and $F^{16}=2I$. Consequently $F^{112}=2I$ on all of $J_C[7]$.
These polynomial identities and both possible Jordan types are
verified by \path{bielliptic_obstructions.py}.

The locally exact theta divisor is Frobenius-stable and symmetric:
if $B^2\simeq\omega_C$, then
$-[R-B]=[\iota R-B]$. A nonzero theta point $a$ of order seven
would therefore force all six nonzero multiples of $a$ into theta,
because $\langle2,-1\rangle=(\mathbf Z/7)^\times$.
This is impossible by the genus-two difference bound.
For a fixed nonzero Jacobian class, there are at most two ordered
representations as $R-S$: equality of two such representations
gives linearly equivalent effective degree-two divisors; if they
are distinct, both are the unique hyperelliptic pencil.
Thus a subset of size $r$ in the intersection of a theta translate
with a cyclic group of order seven satisfies
$r(r-1)\le2(7-1)$. It cannot have $r=6$.
Theta has no nonzero seven-torsion, and the checked locally exact
differential sequence makes every geometric cyclic seventh-degree
Prym ordinary, of two-rank six.

For a nonzero quotient functional $v$ on $J_C[7]$, its period $m$
and a scalar power of Frobenius are exhausted by
\[
 \begin{array}{c|c|c}
 m&d&vF^d/v\\\hline
 3&1&2\\
 21&7&2\\
 24&16&2\\
 168&112&2.
 \end{array}
\]
The first two cases are supported on the $(X-2)$ primary part,
without or with a nonzero Jordan contribution.
The last two include the quadratic primary part.
The certificate checks all $2400$ nonzero vectors for each Jordan
type, as well as the displayed formulas.

The corresponding geometric cover has a pointed model over
$\mathbf F_{4^d}$ with deck Frobenius equal to two.
Its new two-torsion is the sum of the two nontrivial irreducible
$\mathbf F_2[\mathbf Z/7]$-modules, each a one-dimensional
$\mathbf F_8$-space. On each, this Frobenius is semilinear for
squaring, and its cube is the norm to $\mathbf F_2^\times$,
hence the identity. New two-torsion is therefore rational over
$\mathbf F_{4^{3d}}$.
Here $3d=m$ or $2m$. The first field with constant deck group is
$\mathbf F_{4^m}$, and AC.29 shows that rational two-torsion cannot
first appear only in its quadratic extension.
Its relative class number is already even over that first field
and remains even over every extension.
The uniqueness of the pointed descent identifies every eligible
form with this base change. AC.29 also proves the final statement
for power-of-two pointed twists.
\end{proof}


\begin{claim}{AC.31}{Proved counterexample to a scalar-Frobenius reduction}
Let
\[
 C/\mathbf F_{16}:\quad y^2+y=x^5,\qquad P=\infty.
\]
Arithmetic Frobenius on \(J_C\) is the scalar \([-4]\).
Nevertheless there is a constant \(P\)-split cyclic cover of degree
thirteen over \(\mathbf F_{4096}\), its first possible constant-deck
field, whose new factor is ordinary of two-rank twelve and has odd
relative class number. More precisely its twisted polynomials are
\[
 L_P(\chi^i,T)
 =1+\frac{139-13\left(\frac{i}{13}\right)\sqrt{13}}2\,T
       +4096T^2,\qquad 1\le i\le12,
\]
after choosing the sign of \(\sqrt{13}\) consistently with the
quadratic Gauss sum. Consequently
\[
 \frac{P_D(T)}{P_C(T)}
 =\bigl(1+139T+12473T^2+569344T^3+16777216T^4\bigr)^6
\]
and
\[
 \frac{h_{\mathbf F_{4096}}(D)}{h_{\mathbf F_{4096}}(C)}
                         =17359173^6.
\]
Thus scalar Jacobian Frobenius, even at a rational hyperelliptic
ramification point, does not reduce the rational criterion to
vanishing of the geometric new \(p\)-rank.
\end{claim}
\begin{proof}
The exact counts \(\#C(\mathbf F_{16})=33\) and
\(\#C(\mathbf F_{256})=193\) give Frobenius polynomial
\((X+4)^4\). Frobenius acts semisimply, so it equals \([-4]\)
as a Jacobian endomorphism; equivalently apply
\cite[Theorem~2(d), p.~140]{AC_Tate1966}.
The following computation is independent of the assertion that all
thirteenth-degree Pryms of this curve are ordinary: it proves
ordinariness for the displayed cover itself.

Use binary polynomial notation in
\[
 K=\mathbf F_2[t]/(t^{12}+t^3+1),\qquad \#K=4096.
\]
The reduced divisor class
\[
 a=(u,v),\qquad
 u=x^2+3377x+483,\qquad v=3144x+2626
\]
has exact order thirteen. Its sixteenth-power conjugate is
\([-4]a\), and its \(4096\)-power conjugate is \(a\).
These are checked by exact Mumford arithmetic. Since
\((-4)\bmod13\) has order three and \((-4)^2=16\), all nonzero
geometric thirteenth-degree quotient characters have first
constant-deck field \(\mathbf F_{16^3}\). Also
\(4096\equiv1\bmod13\), so the displayed \(a\) has the constant-deck
descent required in AC.27.

For an independent direct character count, write \(u^{13}=M\).
Lift \(v\) to the unique solution of \(w^2+w=x^5\bmod M\) with
\(w\equiv v\bmod u\). Exact linear algebra gives polynomials
\(A,B\), of degrees thirteen and ten, with \(A\) monic and
\[
 A^2+AB+B^2x^5=u^{13},\qquad A+Bw\equiv0\bmod u^{13}.
\]
Their coefficient lists, in increasing degree, are
\[
\begin{aligned}
 A={}&(1334,2010,3185,1642,2095,2911,77,\\
     &\hspace{12mm}567,182,1378,3843,297,2491,1),\\
 B={}&(1250,517,799,137,81,3560,2285,3410,2702,2433,1191).
\end{aligned}
\]
Hence \(F=A+By\) has divisor \(13D_a-26P\), where
\(D_a\) is the effective degree-two divisor representing \(a\).
The norm identity and the vanishing condition distinguish its
two zeroes from their hyperelliptic conjugates.
The cover is the normalization of \(z^{13}=F\).
It is unramified because every valuation of \(F\) is divisible
by thirteen. It is connected because \(a\ne0\), and it is
\(P\)-split because the leading local unit of \(F\) at \(P\)
is one. Its deck group is constant since \(\mu_{13}\subset K\).

Let the primitive field generator be the binary element \(3\);
use its associated primitive thirteenth root to label the
character exponents. At a point where \(F\ne0\), the exponent is
\(\log_3F\bmod13\). At a zero \(R=(r,y)\), the norm identity
shows that the local unit has character exponent
\(-\log_3B(r)\bmod13\): the other factor of \(u^{13}\) contributes
a thirteenth power. Infinity contributes exponent zero.
Enumerating all \(4353\) rational points, with these local units
at the two zeroes, gives
\[
 (N_0,\ldots,N_{12})
  =(399,323,336,323,323,336,336,336,336,323,323,336,323).
\]
Thus the nonzero quadratic residues modulo thirteen have count
\(323\), and the nonresidues have count \(336\).
The quadratic-period identity
\(\sum_{r\ {\rm square}}\zeta_{13}^r=(-1+\sqrt{13})/2\)
gives
\[
 \sum_R\chi([R-P])
 =399+323\frac{-1+\sqrt{13}}2
       +336\frac{-1-\sqrt{13}}2
 =\frac{139-13\sqrt{13}}2.
\]
AC.4, with \(K_C\sim2P\), now gives every displayed
character polynomial and their product, including the exact
relative class number.

As a second check, all twelve Cartier scalars of AC.27 belong to
\(\{72,73\}\subset K\), with six of each value.
Both satisfy \(X^2+X+1=0\). Thus none is zero or one, and
\[
 \overline{P_D(T)/P_C(T)}
   =(1+T+T^2)^6=1+T^2+T^6+T^{10}+T^{12}.
\]
The degree twelve proves geometric new two-rank twelve, while
the value one at \(T=1\) proves the rational two-unit condition.
The program \path{scalar_frobenius_counterexample.py} executes
both the Cartier test and the direct Kummer-point test and
requires their complete polynomials modulo two to agree.

The failure of the proposed general reduction is structural.
On its rank-one \(\mathbf F_{4096}\)-module, the nonconstant
deck Frobenius over \(\mathbf F_{16}\) is semilinear for
\(x\mapsto x^{2^8}\), since \(-4\equiv2^8\bmod13\).
Its cube is a norm to \(\mathbf F_{16}^{\times}\), which need
not be one. Hyperelliptic reflection further restricts the scalar
to \(\mathbf F_4^\times\); it still leaves precisely the two
nontrivial values that occur here. Passing from the witness field
to its cubic constant extension raises these values to one and
destroys the rational criterion, giving an explicit instance of
the unsafe-extension warning in AC.6.
\end{proof}

\begin{claim}{AC.32}{Proved; a simultaneous genus-two counting criterion}
Let \(C/k\) have genus two, with \(k\) a finite field of
characteristic \(p\), and let \(P\in C(k)\).
Let \(\ell\ne p\) be an odd prime such that \(p\) is a primitive
root modulo \(\ell\). Set
\[
 V=J_C(k)/\ell J_C(k),\qquad H=\ell J_C(k),\qquad N=\#C(k),
\]
and assume \(V\ne0\). If either
\[
                  |V|>N\qquad\hbox{or}\qquad N\equiv-1\bmod p,
\]
and
\[
                         (p-1)(N-1)>2(|H|-1),                 \tag{AC.16}
\]
then some surjection \(J_C(k)\to\mathbf Z/\ell\) gives a
\(P\)-split L1 witness over \(k\).

One may replace (AC.16) by the sharper occupancy-dependent
condition
\[
 N_0(N_0-1)+(p-1)(N-N_0)>2(|H|-1),\qquad
 N_0=\#\{R\in C(k):[R-P]\in H\}.
\]
The subgroup \(H\) is the image of multiplication by \(\ell\),
not the much larger kernel of one chosen cyclic character.
No universal assertion about the existence of such a prime and
field is implicit in this conditional criterion.
\end{claim}
\begin{proof}
Suppose every prime cyclic quotient fails the rational criterion.
For each quotient, the nontrivial character evaluations of
\(\Theta=[0]+\sum_{R\in C(k)}[[R-P]\bmod H]\) are permuted
transitively by \(p\)-power Frobenius, because \(p\) is primitive
modulo \(\ell\). Failure means at least one is zero; hence all
are zero. Every nontrivial character of \(V\) factors through
one such quotient, so Fourier inversion in \(\mathbf F_p[V]\)
gives
\[
                 \Theta=c\sum_{b\in V}[b],
                 \qquad c=(N+1)/|V|\in\mathbf F_p.
\]
The denominator is invertible since \(\ell\ne p\).
If \(|V|>N\), some nonzero coset is empty and forces \(c=0\).
The alternative congruence \(N\equiv-1\bmod p\) also forces
\(c=0\) directly. Thus
\[
              N_0\equiv-1\bmod p,\qquad
              N_b\equiv0\bmod p\quad(b\ne0).
\]

For every nonzero Jacobian class \(\delta\), there are at most
two ordered pairs of geometric points \((R,S)\) with
\([R-S]=\delta\).
Indeed \(R-S=R'-S'\) implies
\(R+S'\sim R'+S\). If these effective degree-two divisors are
different, both belong to the unique hyperelliptic pencil;
the only possible second representation is
\((R',S')=(\iota S,\iota R)\).
Consequently the number of ordered distinct pairs in the same
\(H\)-coset is bounded above by \(2(|H|-1)\).
The occupancy congruences give the lower bound
\[
 \sum_bN_b(N_b-1)
   \ge N_0(N_0-1)+(p-1)(N-N_0).
\]
Since \(P\) lies in the zero coset, \(N_0\ge1\); its congruence
improves this to \(N_0\ge p-1\). Therefore
\[
\begin{aligned}
 N_0(N_0-1)+(p-1)(N-N_0)
   &=(p-1)(N-1)+(N_0-1)(N_0-(p-1))\\
   &\ge(p-1)(N-1).
\end{aligned}
\]
Either stated strict inequality contradicts the upper bound.
Thus at least one cyclic quotient is a witness.
The primitive-root hypothesis was used at the first step and
cannot be dropped by this argument.
\end{proof}


\begin{claim}{AC.33}{Proved; a Frobenius obstruction with a split coefficient algebra}
Let \(C/k\) have genus two in characteristic two, with
\(P\in C(k)\), and let \(\ell\ge5\) be an odd prime.
Write \(F\) for arithmetic Frobenius on \(J_C[\ell](\overline k)\).
Suppose that, for some \(d>0\),
\[
 F^d=sI,\qquad s\in\mathbf F_\ell^\times
       \text{ of order }\ell-1,
\]
and that every nonzero quotient functional has exact Frobenius
period \(d(\ell-1)\).
Then no \(P\)-split constant cyclic cover of degree \(\ell\)
satisfies the rational criterion over any finite extension of \(k\).
No assumption that two be primitive modulo \(\ell\) is needed.

In particular this conclusion holds if the characteristic
polynomial of \(F\) is irreducible over \(\mathbf F_\ell\)
and \(q=\#k\) generates \(\mathbf F_\ell^\times\).
\end{claim}
\begin{proof}
The locally exact theta divisor is defined over \(k\).
A nonzero point \(a\) of order \(\ell\) on it would, by the
scalar Frobenius identity, put every nonzero multiple of \(a\)
on theta. In genus two theta is a translate of the curve,
so the difference bound proved in AC.32 would give
\[
                    (\ell-1)(\ell-2)\le2(\ell-1),
\]
which is impossible for \(\ell\ge5\).
Thus theta has no nonzero \(\ell\)-torsion, and every geometric
cyclic \(\ell\)-cover has ordinary new factor, of two-rank
\(\ell-1\).

For any such geometric cover, its kernel is stable under
\(F^d\), and the rational point \(P\) gives its unique pointed
descent to the degree-\(d\) extension \(k_d\).
Its geometric new two-torsion \(W\) is free of rank one over
the finite reduced algebra
\[
              R=\mathbf F_2[X]/\Phi_\ell(X).
\]
Indeed each nontrivial irreducible deck-character component
occurs with multiplicity one when the new factor is ordinary.
The algebra \(R\) may be a product of several fields.
Frobenius over \(k_d\) acts on it by \(X\mapsto X^s\)
(or its inverse with dual character conventions).
This automorphism \(\sigma\) has order \(\ell-1\) and
\[
                              R^\sigma=\mathbf F_2.
\]
To check this last equality, extend scalars to
\(\overline{\mathbf F}_2\), where \(R\) is the algebra of
functions on the nontrivial characters; the primitive scalar
permutes all of them transitively, so only constants remain.
Descent of a constant function to \(R\) requires its value to
lie in \(\mathbf F_2\).

Choose an \(R\)-basis of \(W\). Frobenius is then
\(w\mapsto a\,\sigma(w)\) with \(a\in R^\times\).
Its \((\ell-1)\)-st power is multiplication by
\[
                 a\,\sigma(a)\cdots\sigma^{\ell-2}(a)
                         \in(R^\sigma)^\times=\{1\}.
\]
Therefore all new two-torsion is rational over
\(k_{d(\ell-1)}\). By the exact period hypothesis this is the
first field with constant deck group. The unique pointed
descent identifies every eligible cover over later fields
with its base change, and its rational two-torsion persists.
This proves the exclusion.

For the last assertion, identify the irreducible
four-dimensional Frobenius module with
\(\mathbf F_{\ell^4}\), on which \(F\) multiplies by an
element \(\lambda\).
The symplectic reciprocity involution on its four conjugate
eigenvalues is \(\lambda\mapsto q/\lambda\).
It commutes with their cyclic Galois permutation, is nontrivial,
and has order two. It is therefore the \(\ell^2\)-power map,
giving
\[
                         \lambda^{\ell^2+1}=q.
\]
Since \(q\) generates \(\mathbf F_\ell^\times\),
the cyclic group \(\langle\lambda\rangle\) contains that whole
scalar subgroup. Put \(m=\operatorname{ord}(\lambda)\) and
\(d=m/(\ell-1)\). Then \(\lambda^d\) is a primitive scalar,
while every nonzero vector and quotient functional has exact
period \(m\). The first part applies.
\end{proof}

\paragraph{What the strengthened counts still leave unresolved.}
The preceding statements give a concrete obstruction to treating
the character values in L2 as independent random elements.
Whole primes may fail for representation-theoretic reasons,
before any individual cover is tested. Conversely AC.31 shows
that an ordinary new factor can still be a rational witness even
when Jacobian Frobenius is scalar.

There is also a direct arithmetic limitation on the available
counting criteria. If \(k=\mathbf F_{p^e}\), \(p\) is primitive
modulo \(\ell\), and
\(r=\dim_{\mathbf F_\ell}J_C(k)/\ell J_C(k)>g\), then
\[
                         \ell\le e+1.
\]
Indeed the Frobenius-fixed subspace of \(J_C[\ell]\) is isotropic
for the Weil pairing unless \(q\equiv1\bmod\ell\), and an isotropic
subspace has dimension at most \(g\).
Thus \(r>g\) forces \(\ell\mid q-1\); primitivity gives
\(\ell-1\mid e\).
For large \(e\), a rank greater than \(g\) therefore supplies
only polynomially many cosets, whereas the curve and Jacobian
point counts grow exponentially in \(e\).
In genus two, the threshold in AC.32 consequently still demands
large prime divisors in the lower-rank cases.
No existence theorem for those divisors with the required
residue order has been established here.
The exact Cartier test of AC.27 gives a finite decision for
each proposed character; it does not by itself provide the
uniform torsion-point or Frobenius-orbit avoidance theorem
needed to settle A3 for every curve.


\begin{claim}{AC.34}{Proved; limitation of a trigonal branch-point argument}
Let \(E/L\) be an ordinary elliptic curve in characteristic two,
and let \(f:E\to\mathbf P^1\) be a separable degree-three map
whose ramification consists of three distinct simple wild points
\(R_1,R_2,R_3\). No difference \(R_i-R_j\), \(i\ne j\), is the
nonzero geometric two-torsion class. It is nevertheless possible
for all three differences to be two-primary.

An explicit example over \(\mathbf F_4\), with
\(\alpha^2+\alpha+1=0\) and \(\beta=\alpha^2\), is
\[
 E:\ y^2+xy=x^3+1,\qquad f=(y+\alpha)/x.
\]
The ramification points are
\[
 T=(0,1),\quad R=(\beta,0),\quad R'=(\beta,\beta),
\]
and the pair-difference orders are \(8,8,4\).
The unramified residual points in their degree-three fibres are
\[
 S_\infty=\infty,\qquad S_\beta=(1,1),\qquad S_\alpha=(1,0),
\]
whose pair-difference orders are \(4,4,2\).
The quadratic resolvent is the genus-two curve
\[
 C:\quad y^2+y=z+\frac{\beta}{z^2+z}+\alpha.
\]
Its constant pointed cubic cover has
\(L(\chi,T)=1+3T+4T^2\) for both nontrivial characters and
relative class number \(64\).

For a pointed dihedral cubic cover \(D\to C\), twisting by the
reflection fixing the chosen lift leaves its elliptic quotient
unchanged. Thus this twist cannot reduce the two-primary order
of the elliptic group in the proposed argument.
These statements exhibit a limitation of that route, not a
collision on \(C\).
\end{claim}
\begin{proof}
Let \(\eta\ne0\) be the unique geometric two-torsion point of
\(E\). The locally exact differential line on an ordinary
elliptic curve is the nontrivial two-torsion line bundle
corresponding to \(\eta\).
If \(D_\infty=f^*(\infty)\), the ramification formula gives
\[
 \operatorname{div}(df)
        =2(R_1+R_2+R_3-D_\infty),\qquad
 [R_1+R_2+R_3-D_\infty]=\eta.
\]
For completeness the class of half this exact differential
cannot be zero: otherwise \(df=g^2\omega\) for a nonzero
regular differential \(\omega\); applying Cartier gives
\(0=g\,\mathcal C(\omega)\), contrary to ordinariness.
If \(R_1-R_2=\eta\), cancellation would give
\(D_\infty\sim2R_2+R_3\).
The actual fibre at \(f(R_2)\) is \(2R_2+S_2\).
Cancelling \(2R_2\) and using the injective Abel map of \(E\)
gives \(S_2=R_3\). That fibre would contain both ramification
points with multiplicity at least two, contradicting degree
three. This proves the order-two assertion.

For the example, direct enumeration gives
\(E(\mathbf F_4)\simeq\mathbf Z/8\).
The pole divisor of \(f\) is \(2T+\infty\), so its degree is
three, and with \(\omega=dx/x\) one has
\[
                        df=\frac{x^2+\alpha}{x}\,\omega.
\]
The divisor of this differential and the pole divisor give
different exponent two at exactly \(T,R,R'\).
Their branch values are respectively \(\infty,\beta,\alpha\).
The stated residual points follow by evaluating \(f\);
exact elliptic addition gives the six difference orders in
the claim. In particular the distinction between order two
and a higher two-power is essential.

Writing \(f=t\), so \(y=tx+\alpha\), gives the cubic
\[
                x^3+(t^2+t)x^2+\alpha x+\alpha=0.
\]
For a separable cubic \(X^3+aX^2+bX+c\) in characteristic two,
its quadratic resolvent has Artin--Schreier class
\((b^3+a^3c+c^2)/(ab+c)^2\).
Here this simplifies to
\[
              \beta(t^2+t)+\frac{\beta}{t^2+t+1}.
\]
The change \(w\mapsto w+\alpha t\), followed by \(t=z+\alpha\),
gives the displayed \(C\).
The Galois closure has group \(S_3\), since the cubic is
irreducible over \(L(t)\) by its degree-three function-field
map and the nontrivial quadratic resolvent is connected.
All three branch inertia groups have order two, so the
degree-three map from this closure to \(C\) is unramified.
The rational residual points give rational lifts of the three
branch points of \(C\), making this constant cubic cover
pointed. Alternatively its unique pointed character is
identified directly in \(J_C(\mathbf F_4)\), which has order
twelve. The three rational points of \(C\) are its branch
points and have character occupancies \((3,0,0)\).
AC.4 therefore gives \(1+3T+4T^2\) and relative value
\(8^2=64\). All the finite-field, group, resolvent and
character calculations are executed in
\path{trigonal_two_primary_obstruction.py}.

Finally, if \(D\to C\) has deck group \(\langle\delta\rangle\)
of order three and \(\tau\) is its pointed reflection,
then a \(\tau\)-fixed lift over a branch point maps to the
unramified residual point \(S_i\) of
\(E=D/\langle\tau\rangle\to\mathbf P^1\), not to \(R_i\).
Indeed the total inertia has order two and is already
accounted for in \(D\to E\). The fibre relation is
\[
                     2R_i+S_i\sim D_\infty,\qquad
                     S_i-S_j=-2(R_i-R_j).
\]
Twisting \(D\) by \(\tau\) does not twist \(E\), because
\(\tau\) acts trivially on its quotient. Every other reflection
\(\delta^j\tau\) is conjugate to \(\tau\) by a rational deck
element, so it gives the same elliptic quotient up to
\(L\)-isomorphism when the deck group is constant.
A nontrivial pure deck twist destroys the rational lift over
the prescribed branch point. Thus neither operation supplies
the missing small-two-primary-group hypothesis.
\end{proof}


\begin{claim}{AC.35}{Proved partial reformulation; a finite etale torsor in the scalar case}
Let \(C/\mathbf F_q\) be a genus-two curve in characteristic
two, with a rational base point and the corresponding point-rigidified
Poincare family of degree-zero line bundles, defined over the base field. Suppose arithmetic Frobenius
on \(J_C\) is the scalar \([s]\), where \(s^2=q\).
On a nonempty rational open of \(J_C\), choose a rational basis
\(e_a\) of its one-dimensional nontrivial character cohomology,
and write the \(q\)-power coefficient as
\[
                     F_{\rm abs}^{\log_2q}e_a
                              =A(a)e_{[q]a}.
\]
For the models of AC.27, \(A\) is given explicitly by its
finite Cartier--Mumford product.
On the open \(U\) where this basis and \(A\) are invertible,
\[
                         \mathcal Y:\quad z^{q-1}=A
\]
is a finite etale \(\mu_{q-1}\)-torsor.
Its isomorphism class is independent of the basis on common
domains of definition.

If \(a\in U(\mathbf F_{q^m})\), the linear coherent return
scalar for the pointed character defined by \(a\) is
\[
                    \lambda(a)=
                    N_{\mathbf F_{q^m}/\mathbf F_q} A(a).
\]
This is exactly the arithmetic Frobenius value of
\(\mathcal Y\) at \(a\), with degree \(m\).
If \(a\) has odd prime order \(\ell\) and every nonzero multiple
of \(a\) lies in \(U\), its pointed cyclic cover over
\(\mathbf F_{q^m}\) satisfies (R) exactly when all these
Frobenius values differ from one.

The torsor furnished by \(C:y^2+y=x^5\) over \(\mathbf F_{16}\)
is geometrically nontrivial. This construction supplies a
finite etale object detecting the single-character return
condition, but does not prove that its nontrivial Frobenius
locus contains a whole nonzero prime cyclic subgroup.
\end{claim}
\begin{proof}
The scalar hypothesis gives the identity of Jacobian
endomorphisms
\[
                             [q]=F_J^2.
\]
Since the chosen rational functions are defined over
\(\mathbf F_q\), one has \(A(F_Ja)=A(a)^q\).
For a point \(a\) rational over \(\mathbf F_{q^m}\),
iterating the absolute Frobenius formula gives
\[
\begin{aligned}
 A_m(a)
  &=\prod_{j=0}^{m-1}
      A([q]^j a)^{q^{m-1-j}}\\
  &=A(a)^{\,q^{m-1}(1+q+\cdots+q^{m-1})}\\
  &=N_{\mathbf F_{q^m}/\mathbf F_q}A(a).
\end{aligned}
\]
The last expression lies in \(\mathbf F_q\), so the inverse
coefficient-field conjugation in AC.27 does not alter it.
Also \(F_J^m(a)=a\) implies \([q^m]a=a\).
If \(a\) has order \(n\), then \(n\mid s^m-1\), so
\(\mu_n\subset\mathbf F_{q^m}\).
The point-normalized line bundle therefore defines the
constant pointed cyclic cover over this same field.

For a different rational basis \(e'_a=b(a)e_a\), the coefficient
changes to
\[
 A'(a)=A(a)\frac{b(a)^q}{b([q]a)}
       =A(a)b(a)^{q-q^2}
       =A(a)\bigl(b(a)^{-q}\bigr)^{q-1}.
\]
This is precisely a change of Kummer generator, proving
basis independence. The relevant open is defined over
\(\mathbf F_q\) and is stable under \(F_J\), whose underlying
map is Frobenius; thus the return computation does not
require removing a new infinite collection of divisors.

Over a closed point with residue field \(\mathbf F_{q^m}\),
arithmetic Frobenius sends a root \(z\) to \(z^{q^m}\), so
its ratio is
\[
             z^{q^m-1}
              =A(a)^{(q^m-1)/(q-1)}
              =N_{\mathbf F_{q^m}/\mathbf F_q}A(a).
\]
This proves the exact comparison. For a prime-order
point, all nonzero multiples have the same field of
definition: multiplication by any \(1\le i<\ell\)
has an inverse on that cyclic subgroup.
Applying AC.17 or AC.27 to all its nontrivial characters
gives the stated criterion.

For the example of AC.31, \(q=16\), and the degree-three
closed point \(a\) has return value a nontrivial cube
root in \(\mathbf F_4^\times\).
The program now independently checks that each
\(4096\)-power scalar is the \(4096/16\) norm of its
\(16\)-power Cartier coefficient.
If the \(\mu_{15}\)-torsor were geometrically trivial,
its arithmetic Frobenius at a degree-three point would
be the cube of an element of \(\mu_{15}\), hence have
order dividing five. The observed order three is
impossible. Thus the torsor has nontrivial geometric
monodromy.

For general nonscalar Jacobian Frobenius the corresponding
periodicity condition is \(V_q^m(a)=a\), since on prime-to-two
torsion \(\sigma_{q^m}(a)=q^m a\) is equivalent to that
condition for the \(q\)-Verschiebung.
The simplification \([q]=F_J^2\), essential to the norm
and the Kummer gauge identity above, is then unavailable.
Even in the scalar case, ordinary Chebotarev applied to
\(U\) controls closed points of \(U\); it does not impose
prime order on \(a\) and simultaneous avoidance of one
on every nonzero multiple. Those are the precise additional
assertions still needed for A3.
\end{proof}


\begin{claim}{AC.36}{Proved: an exact nonabelian bounded-exponent example}
Let \(L=\mathbf F_{4096}\), let \(D/L\) be the Fermat curve
\[
                         X^{39}+Y^{39}+Z^{39}=0,
\]
and choose a primitive thirteenth root \(\zeta\in L\).
The automorphisms
\[
 \sigma(X:Y:Z)=(\zeta X:\zeta^3Y:\zeta^9Z),
 \qquad \tau(X:Y:Z)=(Y:Z:X)
\]
generate a constant group \(G\simeq C_{13}\rtimes C_3\)
acting freely on \(D\). For \(C=D/G\), one has
\[
                  g(D)=703,\qquad g(C)=19.
\]
For every odd positive integer \(m\),
\[
 \begin{aligned}
 J_D(\mathbf F_{4096^m})[2^\infty]
                  &\simeq(\mathbf Z/2)^{24},\\
 J_C(\mathbf F_{4096^m})[2^\infty]
                  &\simeq(\mathbf Z/2)^4.
 \end{aligned}
\]
Thus this cover has relative class-number valuation
\[
             v_2\!\left(\frac{h(D)}{h(C)}\right)=20
\]
over every such odd extension: it fails (R).
Nevertheless every \(P\in C(\mathbf F_{4096^m})\), with
\(m\) odd, has a rational lift to this same cover after
an additional extension of degree in \(\{1,3,13\}\),
and the bounded-exponent nonabelian criterion applies there,
since \(39>(2-1)^2\).

For an even extension degree \(m\), the exact exponent of
\(J_D(\mathbf F_{4096^m})[2^\infty]\) is
\(2^{v_2(m)+2}\). In particular it is eight for \(m=2\);
the numerical inequality for this order-39 cover then fails.
The example distinguishes two sufficient tests on a specified
cover; it does not assert that \(C\) has no other (R) witness.
\end{claim}
\begin{proof}
The relation is \(\tau\sigma\tau^{-1}=\sigma^3\).
Every nonidentity element of \(\langle\sigma\rangle\)
has three distinct projective eigenvalues, and its fixed
points are the coordinate vertices, none on \(D\).
Every remaining nonidentity element of \(G\) has order
three and is conjugate to \(\tau\) or \(\tau^2\).
Their projective fixed points have coordinates
\((1:\omega:\omega^2)\), with \(\omega^3=1\),
and at such a point the Fermat expression is
\(1+1+1=1\). Thus the action is free.
The plane genus formula and the etale Riemann--Hurwitz
formula give the two stated genera.

Here are the complete spectral data, so that the claim
about the integral two-primary exponent is independently
checkable. Put \(q=4096\), choose a multiplicative
character \(\chi:L^\times\to\mu_{39}\), and set
\[
 J(a,b)=\sum_{x\in L\setminus\{0,1\}}
                   \chi(x)^a\chi(1-x)^b,
 \qquad \alpha_{a,b}=-J(a,b).
\]
The indices are \(1\le a,b\le38\), \(a+b\not\equiv0\pmod{39}\).
For the diagonal group \(\mu_{39}^2\), each such character
occurs exactly once in \(H^1\), and its Frobenius scalar is
\(\alpha_{a,b}\). This assertion, including its sign and its
validity for both prime-to-two etale and crystalline cohomology,
is exactly \cite[Corollary~2.4, p.~181]{AC_Katz1981}.
The integral crystalline version is explained on p.~183;
Proposition~4.2 on pp.~185--186 also transfers the character
polynomial identities to the Jacobian endomorphism ring.

The accompanying program evaluates all \(1406\) sums in
\(\mathbf Z[\zeta_{39}]\), using the binary field
\[
              L=\mathbf F_2[t]/(t^{12}+t^3+1)
\]
and the primitive multiplicative generator encoded by the
integer \(3\). It checks each equality
\(\alpha\bar\alpha=4096\), each invariance under
\(\zeta_{39}\mapsto\zeta_{39}^2\), and the following table.
The conductor is \(39/\gcd(a,b,39)\); the second column
is \(\alpha+\bar\alpha\), including twice the scalar in
the rational cases.
\[
\begin{array}{r|r|r}
 \text{conductor}&\alpha+\bar\alpha&\text{character pairs}\\ \hline
 3&128&2\\
 13&-128&132\\
 39&-128&360\\
 39&-80&216\\
 39&-28&72\\
 39&28&432\\
 39&80&144\\
 39&115&48
\end{array}
\]
In particular the possible trace \(-115\) does \emph{not}
occur. As an independent sign check, direct enumeration
of projective points gives \(\#D(L)=30537\), and the
sum of the displayed eigenvalues is \(-26440=q+1-30537\).
Consequently the zeta numerator is
\[
\begin{aligned}
 P_D(T)={}&(1-64T)^2(1+64T)^{492}\\
 &\cdot(1+80T+4096T^2)^{108}
        (1+28T+4096T^2)^{36}\\
 &\cdot(1-28T+4096T^2)^{216}
        (1-80T+4096T^2)^{72}\\
 &\cdot(1-115T+4096T^2)^{24}.
\end{aligned}
\]

It is essential to prove an exponent bound on the actual
Tate lattice, rather than infer one from an unspecified
isogeny decomposition. Let \(F\) denote the \(L\)-Frobenius
endomorphism of \(J_D\). It is central in the semisimple
algebra \(\operatorname{End}_L^0(J_D)\); hence its minimal
polynomial over \(\mathbf Q\) is square-free. By the
faithful prime-to-two Tate realization, its roots occur
among the reciprocal roots of the displayed zeta numerator. The action
of \(F\) on \(V_2^{\rm et}(J_D)\) remains semisimple,
and all its eigenvalues are two-adic units, since \(F\)
is an automorphism of its integral etale Tate lattice.

There is only one unit root among the corresponding
Frobenius polynomials: the root \(\beta\in\mathbf Z_2^\times\) of
\[
                       X^2-115X+4096=0.
\]
The other root \(\gamma\) has valuation twelve.
For the traces \(\pm80\) and \(\pm28\), both roots
have positive valuation, as do the rational roots
\(\pm64\). Thus semisimplicity gives
\[
               F=\beta\,\operatorname{id}
               \quad\text{on }V_2^{\rm et}(J_D),
\]
and therefore also on the actual integral Tate lattice.
This argument excludes any hidden lattice extension.
Since
\[
 (1-\beta)(1-\gamma)=4097-115=3982=2\cdot1991,
\]
one has \(v_2(\beta-1)=1\). The rational two-primary
subgroup is therefore elementary abelian. Its order is
\(2^{24}\), directly from \(P_D(1)\), proving its exact
structure without needing to infer it merely from
Newton slopes.

For completeness, the quotient calculation is also made
at the level of individual characters. Write
\(c=-a-b\pmod{39}\). The \(\sigma\)-invariant characters
are exactly those with
\[
                    a+3b+9c\equiv5a+7b\equiv0\pmod{13}.
\]
There are \(110\) of them. Under the cycle
\((a,b,c)\mapsto(b,c,a)\), two are fixed,
namely \((13,13,13)\) and \((26,26,26)\), and the
other \(108\) form \(36\) orbits of length three.
Each length-three orbit contributes one invariant line,
with its common Jacobi eigenvalue. The two fixed lines
factor through the Fermat cubic. The coordinate cycle
acts freely on that genus-one curve and hence by a
translation, so it acts trivially on these two lines.
The invariant eigenvalues are \(64\) twice, \(-64\)
twenty-eight times, and the two roots of
\(X^2-115X+4096\) four times each. Thus
\[
        P_C(T)=(1-64T)^2(1+64T)^{28}
                           (1-115T+4096T^2)^4.
\]
The same scalar argument, or the pullback embedding
on two-adic Tate modules for the odd-degree cover,
gives the elementary abelian group of order \(2^4\).
In particular \(f_2(D)=24\), \(f_2(C)=4\), and
\(v_2(P_D(1)/P_C(1))=24-4=20\).

The field claims now follow with no change of cover.
The root \(\beta\) is congruent to \(115\pmod{4096}\),
since \(\beta=115-\gamma\). Therefore
\[
                   v_2(\beta-1)=1,\qquad
                   v_2(\beta+1)=2.
\]
For odd \(m\), the factor
\(1+\beta+\cdots+\beta^{m-1}\) is a unit, so
\(v_2(\beta^m-1)=1\). For even \(m\), the elementary
two-adic lifting formula gives
\[
                  v_2(\beta^m-1)=v_2(m)+2.
\]
As \(F\) is scalar on the lattice, these are exact
exponents of the finite rational groups, not only bounds.

For any rational point of \(C\) over an odd-degree
extension of \(L\), its geometric fibre is a \(G\)-torsor.
Frobenius acts there as multiplication by some element
of \(G\), up to conjugacy; its order is one, three or
thirteen. Thus an additional extension of that odd
degree makes the fibre split. The exponent stays two.
The nonabelian annihilator argument of AG.5, with
\(a=2\cdot1\in\mathbf Z[G]\), now excludes a collision
because its degree condition is \(39>(2-1)^2\).
For even extension degree two the exponent is eight
and \(39\le(8-1)^2=49\), which explains the stated
limitation of this particular numerical test.
\end{proof}

\paragraph{Scope of the example.}
There is a separate geometric reason that this particular
quotient does not provide a new unresolved instance of A1.
The two power maps from the Fermat curve of degree \(39\)
to the Fermat curves of degrees three and thirteen are
jointly injective on geometric points: in any common
nonzero coordinate chart, equality of both images forces
ratios whose third and thirteenth powers are both one,
hence forces equal ratios. Both target Jacobians have
two-rank zero, as is already shown by the conductor-three
and conductor-thirteen entries of the Jacobi table.
A two-primary point difference on \(D\) therefore maps
to zero in both Jacobians; the Abel--Jacobi embeddings
and joint injectivity force the two points to coincide.
The collision-lifting statement AC.1 transfers this
exclusion to \(C\), over every finite extension. The
point of AC.36 is the exact nonabelian exponent test,
including a cover for which (R) fails; it is not a claim
that all-extension injectivity of this example required
that numerical test.

\begin{claim}{AC.37}{Proved: a free order-117 group and its precise limitation}
On the Fermat curve \(D/\mathbf F_{4096}\) of AC.36,
choose a primitive thirty-ninth root \(\xi\) and put
\[
 s(X:Y:Z)=(X:\xi Y:\xi^{17}Z),
 \qquad \tau(X:Y:Z)=(Y:Z:X).
\]
The constant nonabelian group
\[
                 G'=\langle s,\tau\rangle
                    \simeq C_{39}\rtimes_{16}C_3
\]
has order \(117\) and acts freely. Its quotient \(C_7=D/G'\)
has genus seven and zeta numerator
\[
 P_{C_7}(T)=(1-64T)^2(1+64T)^4
                         (1-115T+4096T^2)^4.
\]
Thus \(J_{C_7}(\mathbf F_{4096})[2^\infty]
\simeq(\mathbf Z/2)^4\), and \(D/C_7\) again has relative
class-number valuation twenty.
The bounded-exponent criterion for this cover applies
after splitting the pointed fibre over every extension
of \(\mathbf F_{4096}\) of odd or twice-odd degree:
the relevant exponents are respectively two and eight,
and \(117>(8-1)^2\).

There is also a more direct (R) witness for this particular
quotient. The subgroup \(\langle s^3,\tau\rangle\) is
normal of index three and is the order-39 group of AC.36,
after choosing its generator of order thirteen.
Writing \(C_{19}=D/\langle s^3,\tau\rangle\), the cover
\(C_{19}\to C_7\) is etale cyclic of degree three, with
relative zeta numerator and relative class number
\[
                  (1+64T)^{24},\qquad 65^{24}.
\]
Its new two-rank is zero, so (R) persists over every
extension, and every prescribed base point becomes
split after an additional extension of degree at most
three. Accordingly, the smaller genus of this quotient
does not enlarge the set of curves reached by (R).
\end{claim}
\begin{proof}
Projectively,
\[
       \tau s\tau^{-1}=s^{16},\qquad
       16^3\equiv1\pmod{39},\qquad
       1+16+16^2\equiv0\pmod{39}.
\]
The diagonal generator has order \(39\), and no nonidentity
power of \(\tau\) is diagonal, giving the stated order
and semidirect product. The pairwise differences of the
three diagonal exponents \(0,1,17\) are \(1,16,17\),
all units modulo \(39\). Hence a nonidentity diagonal
element has distinct eigenvalues and fixes only the
coordinate vertices.

For an element \(s^a\tau^b\), with \(b=1\) or \(2\),
use its displayed monomial matrix. Its cube is
\(\xi^{18a}I\). If \(\lambda\) is an eigenvalue,
\[
              \lambda^3=\xi^{18a},\qquad
              \lambda^{39}=\xi^{234a}=1.
\]
All eigenvector coordinates are nonzero; their ratios
are products of powers of \(\lambda\) and \(\xi\),
so their thirty-ninth powers are equal. The Fermat
expression at an eigenpoint is therefore a nonzero
scalar times \(1+1+1=1\). There are no fixed points.
Riemann--Hurwitz gives
\(g(C_7)=1+(703-1)/117=7\).

For the Jacobi character indexed by \((a,b,c)\),
\(c=-a-b\), the diagonal invariance condition is
\[
                   b+17c\equiv22a+23b\equiv0\pmod{39}.
\]
The augmented exact program finds \(38\) characters:
the same two cycle-fixed characters as in AC.36 and
twelve cycles of length three. The quotient eigenvalues
are \(64\) twice, \(-64\) four times, and the two
roots of \(X^2-115X+4096\) four times each.
This proves the displayed numerator. The actual
integral exponent follows from the scalar unit-root
argument of AC.36.

Every element of \(G'\) has order in
\(\{1,3,13,39\}\). A rational fibre therefore splits
after an odd extension of one of those degrees.
Such an extension does not change the two-adic valuation
of the original constant-field degree. Hence the exact
exponent remains two for odd degrees and eight for
twice-odd degrees, as stated.

Finally, the map \(G'\to C_3\) sending \(s\) to a
generator and \(\tau\) to zero is well-defined because
\(16\equiv1\pmod3\). Its kernel is
\(\langle s^3,\tau\rangle\); the order-thirteen
generator \(s^6\), after multiplying coordinates by
a common scalar, is the diagonal generator in AC.36.
Dividing the two explicitly proved zeta numerators gives
\((1+64T)^{24}\). Its roots all have positive two-adic
valuation; alternatively its value over every extension
is an odd power of \(1-(-64)^m\). Thus (R) persists
under the extension used to split any given fibre.
\end{proof}

\begin{claim}{AC.38}{Proved: when scalar Frobenius forces the arithmetic test to be geometric}
Let \(C/k\) be a genus-two curve in characteristic two,
let \(P\in C(k)\), and suppose Jacobian Frobenius is the
scalar \([s]\), where \(s=\varepsilon2^r\),
\(\varepsilon\in\{1,-1\}\), and \(\#k=s^2\).
Let \(\ell\) be an odd prime and put
\[
 f=\operatorname{ord}_\ell(2),\qquad
 m=\operatorname{ord}_\ell(s),\qquad
 d=\frac{f}{\gcd(f,m)}.
\]
Every geometric cyclic degree-\(\ell\) cover has a
\(P\)-pointed descent to \(k\), possibly with nonconstant
deck group, and its first constant-deck field has degree
\(m\). Suppose
\[
                  d=1,
        \quad\text{equivalently}\quad
          \langle2\rangle\subseteq\langle s\rangle
                         \text{ in }\mathbf F_\ell^\times.
\]
Then all its geometric new two-torsion points are rational over
that first constant-deck field. In particular this
cover satisfies (R) over a field with constant deck
group if and only if its geometric new two-rank is zero.
This equivalence holds individually for every cyclic
cover, and persists over all further constant fields.
No assumption that two is primitive modulo \(\ell\)
is required.

The condition has the following explicit forms:
\[
\begin{array}{c|c}
 s&\text{condition for }d=1\\ \hline
 +2^r&\gcd(r,f)=1\\
 -2^r,\ f\text{ even}&\gcd(r+f/2,f)=1\\
 -2^r,\ f\text{ odd}&\gcd(r,f)=1.
\end{array}
\]
When \(d>1\), the same argument yields a norm in
\(\mathbf F_{2^d}^\times\), which need not be one.
It therefore gives no converse obstruction in that case.
\end{claim}
\begin{proof}
The scalar action preserves every geometric degree-\(\ell\)
cyclic quotient of the prime-to-two geometric fundamental
group. The rational point supplies its pointed descent,
with arithmetic Frobenius acting on the deck quotient
by the scalar \(s\), up to the harmless simultaneous
inverse convention. Thus the deck group first becomes
constant after degree \(m\).

Consider the new part of geometric etale cohomology
\[
 W=H^1_{\rm et}(\overline D,\mathbf F_2)_{\rm new},
 \qquad R=\mathbf F_2[X]/\Phi_\ell(X).
\]
The averaging idempotent is defined because \(\ell\)
is odd, so \(W\) is an \(R\)-module.
The Artin--Schreier comparison identifies its scalar
extension with the Frobenius-stable part of the
corresponding coherent cohomology. Every nontrivial
character of the geometric cyclic cover occurs in
\(H^1(\overline D,\mathcal O_{\overline D})\) with
dimension \(g(C)-1=1\). Consequently \(W\) has rank
zero or one over each simple field factor of \(R\).
This is the only step using genus two; compare the
functorial stable-part comparison in AC.17.

Arithmetic Frobenius \(F\) is an automorphism of \(W\)
and is semilinear for
\(\sigma_s:R\to R\), \(X\mapsto X^s\).
If \(W\ne0\), its support is a union of \(\sigma_s\)-orbits
of simple factors. Thus \(W\) is free of rank one over a
\(\sigma_s\)-stable direct factor \(R'\) of \(R\).
Choose a basis on that factor; then
\[
                 F=c\sigma_s\quad(c\in(R')^\times),
 \qquad
                 F^m=\prod_{j=0}^{m-1}\sigma_s^j(c).
\]
The last product belongs to \(((R')^{\sigma_s})^\times\).

Here is the exact fixed-ring calculation. The primitive
\(\ell\)-th roots are indexed by
\(\mathbf F_\ell^\times\); their two-Frobenius orbits
have length \(f\), while multiplication by \(s\) has
order \(m\). As both acting subgroups lie in the same
cyclic group, their intersection has order \(\gcd(f,m)\).
It follows that
\[
 R^{\sigma_s}\simeq
 \left(\mathbf F_{2^{\,f/\gcd(f,m)}}\right)^{
                    (\ell-1)/\operatorname{lcm}(f,m)}.
\]
For \(d=1\) its unit group, and that of every direct
factor such as \((R')^{\sigma_s}\), is trivial.
Therefore \(F^m=1\) on \(W\).
The same holds on the dual geometric new two-torsion
of the Jacobian. If \(W\ne0\), criterion (R) fails at
the first constant field and at every later one. If
\(W=0\), the geometric new two-primary group is already
zero, so (R) holds over every extension. This proves the
equivalence, including its field quantifier.

For \(s=2^r\), one has \(m=f/\gcd(r,f)\), giving
\(d=\gcd(r,f)\). If \(s=-2^r\) and \(f\) is even,
then \(-1=2^{f/2}\pmod\ell\), so replace \(r\) by
\(r+f/2\). If \(f\) is odd, inversion is outside
\(\langle2\rangle\), and
\(m=2f/\gcd(r,f)\); because \(f\) is odd this gives
\(d=\gcd(r,f)\) again.
\end{proof}

\paragraph{Relation to the ordinary-new counterexample.}
For \(C:y^2+y=x^5\) over \(\mathbf F_{16}\), one has
\(s=-4\) and \(r=2\). At \(\ell=7\), \(f=3\) and
\(d=1\); at \(\ell=11\), \(f=10\) and
\(d=\gcd(2+5,10)=1\). For those primes the arithmetic
criterion cannot improve on the geometric new-rank-zero
criterion. At \(\ell=13\), however, \(f=12\) and
\(d=\gcd(2+6,12)=4\). The semilinear norm lies in
\(\mathbf F_{16}^\times\), and AC.31 exhibits the
nontrivial values of order three which give an ordinary-new
rational (R) witness. Thus AC.38 explains exactly why
the unrestricted scalar-Frobenius assertion is false,
while recovering a substantial range in which it is true.

\begin{claim}{AC.39}{Proved: exact finite-level meaning of hereditary relative tests}
Let \(G\) be a finite solvable group of order prime to \(p\),
and let \(V\) be a finite \(p\)-primary abelian group with a
\(G\)-action. For a subgroup \(H\), write
\[
                 e_H=\frac1{|H|}\sum_{h\in H}h
                      \quad\text{on }V.
\]
The following conditions are equivalent:
\begin{enumerate}
\item For every pair of normal subgroups \(H<K\) of \(G\),
      one has \(V^H\supsetneq V^K\).
\item There is \(v\in V\) such that
      \(\operatorname{Stab}_G(e_Hv)=H\) for every normal
      subgroup \(H\) of \(G\).
\end{enumerate}

In particular, let \(D/k\to C/k\) be a constant solvable
prime-to-\(p\) Galois etale cover, with a rational chosen
point over \(P\), and put \(V=J_D(k)[p^\infty]\).
All relative class-number tests for
\[
                    D/H\longrightarrow D/K,
                    \qquad H<K\text{ normal in }G,
\]
fail if and only if there is a rational \(p\)-primary
class on \(D\) whose norms have free deck orbits on
every normal quotient of this particular cover.
This is a finite-level statement about classes, with
no assertion that they are differences of points.
\end{claim}
\begin{proof}
Necessity is immediate: if \(V^H=V^K\), then \(e_Hv\)
is \(K\)-fixed, so cannot have stabilizer exactly \(H\).
For sufficiency, induct on \(|G|\). The trivial group
causes no difficulty. Choose a nontrivial abelian normal
subgroup \(N\) of the solvable group \(G\), and split
\[
                  V=U\oplus W,\qquad
                  U=e_NV=V^N,\quad W=(1-e_N)V.
\]
Both summands are \(G\)-stable. The strict-invariant
condition descends to the \(G/N\)-module \(U\), so induction
gives \(u\in U\) with the required property for all normal
quotients of \(G/N\). For any normal \(H\) of \(G\),
\(e_Hu=e_{HN}u\), whence
\[
                     \operatorname{Stab}_G(e_Hu)=HN.
\]

We choose \(w\in W\) simultaneously for all the normal
subgroups \(H\). The algebra \(\mathbf Z_p[N]\) is finite
etale and commutative; it is a product of unramified
discrete valuation rings, and \(N\) acts through a
prime-to-\(p\) character on each factor. The projectors
\(e_H\) are central in \(\mathbf Z_p[G]\). They commute
with this algebra and with each other, so decompose
\(W\) further into their finitely many joint image/kernel
summands. Choose one nonzero element in every nonzero
joint summand and add these elements to obtain \(w\).

If a prime-to-\(p\) root of unity is not one, its
difference from one is a unit in the relevant
unramified valuation ring. Thus every nonzero element
of a character summand has precisely that character's
kernel as its stabilizer in \(N\), even though the
module need not be free over that ring. Our choice
therefore ensures, simultaneously for all normal \(H\),
\[
 \operatorname{Stab}_N(e_Hw)
       =\ker\bigl(N\longrightarrow
                         \operatorname{Aut}(W^H)\bigr).
\]
Denote this kernel by \(K_H\). It is normal in \(G\),
because \(N\) and \(H\) are normal and \(W^H\) is
\(G\)-stable. It contains \(H\cap N\). If this
containment were strict, then \(HK_H>H\); both would
be normal in \(G\), while \(K_H\) acts trivially on
both \(U^H\) and \(W^H\). This would give
\(V^H=V^{HK_H}\), contrary to the hypothesis.
Consequently \(K_H=H\cap N\).

Set \(v=u+w\). Since \(U\) and \(W\) are \(G\)-stable,
an element stabilizing \(e_Hv\) must stabilize each
component separately. From the first component it
belongs to \(HN\). Modulo \(H\) it is therefore
represented by some \(n\in N\); from the second
component this \(n\) lies in \(H\cap N\).
The stabilizer is exactly \(H\), as required.

For the geometric application, the prime-to-\(p\)
pullback and norm identities identify
\(J_{D/H}(k)[p^\infty]\) with \(V^H\). Therefore
\[
 v_p\!\left(\frac{h(D/H)}{h(D/K)}\right)>0
                  \quad\Longleftrightarrow\quad
                  V^H\supsetneq V^K.
\]
The pullback of the actual norm of a class \(v\)
is \(|H|e_Hv\). Multiplication by \(|H|\) is
invertible on \(V\), so it does not change its
stabilizer. The abstract equivalence gives exactly
the asserted finite-level application.
\end{proof}

\paragraph{Why tests only over the original base are weaker.}
Let \(p=2\), \(G=C_7\rtimes C_3\), and let \(V\) be
the additive group of \(\mathbf F_4\), with the action
inflated from \(G^{\rm ab}=C_3\) by multiplication by
a primitive cube root. The proper normal subgroups
are \(1\) and \(C_7\), and
\[
                     V^G=0,\qquad
                     V^1=V^{C_7}=V\ne0.
\]
Thus every nontrivial normal quotient over the original
base has a nonzero new module, and all norm maps have
their usual prime-to-two splittings. Nevertheless every
vector is fixed by \(C_7\), so no free orbit exists.
The omitted relative test is exposed by
\(V^1=V^{C_7}\). This is an algebraic countermodel to
the inference from those module conditions alone;
it is not asserted to be the full Jacobian module of
a specified curve or a counterexample to A1 or A10.

\paragraph{Remaining inverse-system issue.}
AC.39 treats a finite cover with constant deck group.
At a fixed finite field, such covers need not be
cofinal in the geometric prime-to-\(p\) tower.
For nonconstant deck groups, the rational subgroup
\(J_D(k)[p^\infty]\) need not be \(G\)-stable, so the
commuting-projector argument cannot simply be repeated
on that subgroup. Making the deck group constant
produces regular classes over an extension, but does
not make them fixed by the original Frobenius.
Moreover regularity on a cofinal selection of levels
does not imply regularity after norms to every
geometric normal quotient. A semilinear descent and
compatibility statement resolving these two issues
is still required before AC.39 could imply a
Frobenius-fixed element of the full ghost module.

\begin{claim}{AC.40}{Proved: semilinear simultaneous regularity from all relative tests}
Let \(G\) be a finite solvable group of order prime to \(p\),
let \(V\) be a finite-dimensional \(\mathbf F_p\)-vector space
with a \(G\)-action, and let \(F\in\operatorname{GL}(V)\) and
\(\phi\in\operatorname{Aut}(G)\) satisfy
\[
                         F(gv)=\phi(g)F(v).
\]
Assume that for every positive integer \(r\) and every
\(\phi^r\)-stable pair of subgroups \(A\triangleleft B\)
with \(B\triangleleft G\) and \(B/A\) cyclic of prime order, one has
\[
                  (V^A)^{F^r}\supsetneq(V^B)^{F^r}.
                  \tag{AC.40.1}
\]
Then there is \(v\in V^F\) such that
\[
                    \operatorname{Stab}_G(e_Hv)=H
                    \quad\text{for every normal }H\triangleleft G.
                    \tag{AC.40.2}
\]
The subgroups \(H\) in this conclusion need not be
\(\phi\)-stable.

For a pointed geometrically Galois odd prime-to-\(p\)
etale cover \(D/k\to C/k\), take
\(V=J_D[p](\overline k)\), with its actual arithmetic
Frobenius. Failure of every relative cyclic-prime
class-number test over intermediate pointed bases which are
geometrically normal over \(C\), and over every finite
constant extension, implies (AC.40.1).
Consequently it supplies the class in (AC.40.2).
The rational class need not be a point difference.
\end{claim}
\begin{proof}
We induct on \(|G|\). Choose a minimal nontrivial subgroup
\(N\subseteq G\) which is normal in \(G\) and stable under
\(\phi\). It is elementary abelian of order a power of
some prime \(\ell\ne p\). Indeed its derived subgroup,
its Sylow subgroups, and its subgroup of \(\ell\)-th
powers would otherwise give smaller characteristic
subgroups; solvability rules out the nonabelian
characteristically simple case.

The averaging projector gives a \(G,F\)-stable splitting
\[
                         V=U\oplus W,\qquad
                         U=V^N,\quad W=(1-e_N)V.
\]
The hypothesis descends to the \(G/N\)-module \(U\),
so induction gives \(u\in U^F\) satisfying (AC.40.2)
for that quotient. For every normal \(H\) in \(G\),
\[
                         \operatorname{Stab}_G(e_Hu)=HN.
                         \tag{AC.40.3}
\]

Decompose \(W\) by the commuting projectors \(e_H\),
for all normal subgroups \(H\) of \(G\), and by the
simple factors of the commutative semisimple algebra
\(\mathbf F_p[N]\). The former joint blocks are
indexed by kernels \(L\triangleleft G\) of irreducible
\(G\)-constituents: on a block with kernel \(L\),
\[
                         e_H=1\quad\Longleftrightarrow\quad H\subseteq L.
\]
This follows from Maschke's theorem, since for normal
\(H\) the \(H\)-invariants of an irreducible constituent
are either the whole constituent or zero. We group
together all constituents with the same kernel; no
choice of a splitting field is needed for that grouping.
The further \(N\)-character factors correspond to
\(p\)-power orbits of nontrivial characters of \(N\).
All characters in such a factor have the same kernel.

The operator \(F\) permutes these refined blocks.
For each orbit of blocks whose direct sum has a
nonzero \(F\)-fixed space, choose a nonzero fixed
vector on that orbit, and add these choices to obtain
\(w\in W^F\). Each chosen vector has a nonzero component
in every block of its orbit. On an \(N\)-character
factor, a nontrivial character value minus one is
invertible; thus every nonzero component has exactly
the corresponding character kernel as its stabilizer
in \(N\). Consequently, simultaneously for all normal
\(H\), our \(w\) realizes the pointwise kernel:
\[
 \operatorname{Stab}_N(e_Hw)=
 \ker\bigl(N\text{ acting on the set }e_H(W^F)\bigr).
                         \tag{AC.40.4}
\]
Here \(e_H(W^F)\) need not be a \(G\)-submodule or
even equal \((W^H)^{F^r}\) for a field of definition
of \(H\). We do not make either of these replacements.

We prove that the kernel in (AC.40.4) is \(H\cap N\).
Suppose that it contains \(n\in N\setminus H\).
If \(n\) acted trivially on \(V^H\), so would all its
\(G\)-conjugates, because \(V^H\) is \(G\)-stable. Put
\[
 B=H\langle gng^{-1}:g\in G\rangle.
\]
Then \(B\triangleleft G\), \(B/H\) is nontrivial
elementary abelian of exponent \(\ell\), and
\(V^H=V^B\). Choose a hyperplane \(A/H\) in \(B/H\).
It gives \(A\triangleleft B\), \(B/A\simeq C_\ell\),
and \(V^A=V^B\), contradicting (AC.40.1) after a
common power making \(F\) and \(\phi\) the identity.
Thus \(n\) does not act trivially on \(V^H\).
Since \(N\) acts trivially on \(U\),
there is an \(N\)-character \(\chi\) in \(W^H\)
with \(\chi(n)\ne1\).

Among all normal constituent kernels \(L\supseteq H\)
supporting a nonzero character in the line
\[
                 \{\chi^a: a\in\mathbf F_\ell^\times\},
\]
choose one maximal under inclusion. Put
\(K=\ker\chi\subset N\). A surviving constituent with
kernel \(L\) shows \(L\cap N\subseteq K\), so
\[
                         LK\triangleleft LN,\qquad
                         LN/LK\simeq C_\ell.
\]
Modulo \(L\), the first subgroup is
\(K/(L\cap N)\) inside the abelian group
\(N/(L\cap N)\), which proves the asserted normality.
Consider the relative module
\[
                         V^{LK}/V^{LN}.
                         \tag{AC.40.5}
\]
Its constituents have a normal \(G\)-kernel containing
\(L\), and a nontrivial \(N\)-character in the displayed
line. Maximality therefore forces their \(G\)-kernel
to be exactly \(L\). Thus (AC.40.5) consists precisely
of the \(L\)-kernel blocks with nontrivial \(N\)-character
kernel \(K\).

No full \(F\)-orbit meeting this relative module can
have an \(F\)-fixed vector. If it did, our choice of
\(w\) would put a nonzero component in its selected
block. Since \(H\subseteq L\), that component survives
\(e_H\); every nonzero multiple of \(\chi\) detects
\(n\). This contradicts \(n\in\operatorname{Stab}_N(e_Hw)\).

Let \(r\) be the least positive integer for which
\(\phi^r(L)=L\) and \(\phi^r(K)=K\). The pair \((L,K)\)
is determined by each refined block in (AC.40.5).
Hence \(r\) divides the length \(t\) of its full
\(F\)-orbit of blocks. The part of that orbit appearing
in (AC.40.5) is exactly one cycle under \(F^r\).
A fixed vector of \(F^r\) on this cycle is equivalent
to a fixed vector of \(F^t\) on one block, and hence
to a fixed vector of \(F\) on the full orbit.
There is none by the preceding paragraph.

The averaging splitting of (AC.40.5) is \(F^r\)-stable,
so we conclude
\[
                         (V^{LK})^{F^r}=(V^{LN})^{F^r}.
\]
This contradicts (AC.40.1) for the cyclic-prime pair
\(LK\triangleleft LN\). Therefore
\(\operatorname{Stab}_N(e_Hw)=H\cap N\) for every normal
\(H\). Combining this with (AC.40.3) in the
\(G\)-stable direct sum \(U\oplus W\) proves that
\(v=u+w\) satisfies (AC.40.2), exactly as in AC.39.

For the geometric application, odd-order groups are solvable
by the verified Feit--Thompson theorem used in AC.3.
A \(\phi^r\)-stable
subgroup defines its pointed intermediate curve over
\(k_r\). Prime-to-\(p\) norm and pullback identify
the relevant \(p\)-torsion invariants with
\(V^A,V^B\); their strictness is exactly the presence
of rational new \(p\)-torsion, hence failure of (R)
for \(D/A\to D/B\) over \(k_r\). The deck group of
this relative cyclic-prime cover may be nonconstant.
If (R) succeeded, AC.20 would make that cyclic group
constant after its exact Frobenius orbit length,
raising the relative class number to that power and
preserving its \(p\)-unit property. Thus failure of
all constant-deck relative tests over every extension
indeed implies the stated hypothesis.
\end{proof}

\paragraph{The normal-level distinction remains necessary.}
The conclusion concerns all normal quotients of \(D/C\).
The relative cyclic cover isolated in the proof can
have source \(D/(LK)\) which is not normal over \(C\),
although its base \(D/(LN)\) is normal over \(C\).
Accordingly the compactness consequence for the
normal-level ghost module is a one-way implication
from hereditary failure of all relative tests.
The reverse assertion needs freeness at the nonnormal
source as well; it does not follow merely from
regularity at every normal quotient.


\subsection{The hyperelliptic refinement of the scalar obstruction}

\begin{claim}{AC.41}{Proved; an all-extension obstruction using the pointed involution}
Let $C/k$ be a genus-two curve over a finite field of
characteristic two, and let $P\in C(k)$ be a Weierstrass
point. Fix a nontrivial geometric connected cyclic etale
cover of odd degree $n$, with a chosen point above $P$.
Let $k_d/k$ be the first constant extension over which
its geometric isomorphism class is stable, and use its
unique pointed descent to $k_d$. Suppose Frobenius on
its geometric cyclic deck group is multiplication by
$a\in(\mathbf Z/n)^\times$, of order $m$. Put
\[
 U=(\mathbf Z/n)^\times,\qquad
 T=\langle2\rangle\subset U,\qquad
 H=\langle a,-1\rangle\subset U.
\]
If $T\subset H$, every geometric two-torsion point of
the primitive factor is rational over the first
constant-deck field $k_{dm}$. Consequently, for this
pointed cover, the primitive class number is odd over
some constant-deck field if and only if its geometric
primitive two-rank is zero. When $n=\ell$ is prime,
this is the whole relative criterion \textup{(R)}.

More generally, the possible primitive Frobenius
multipliers over $k_{dm}$ belong to the units of a
direct factor of
\[
 \left(\mathbf F_{2^{\delta}}\right)^c,
 \qquad
 \delta=\frac{|T|}{|T\cap H|},\qquad
 c=\frac{|U|}{|TH|}.                              \tag{AC.41.1}
\]
Thus the preceding sufficient condition is exactly
$\delta=1$ for this fixed-ring argument.

In particular, suppose the Jacobian Frobenius of $C/k$
is the scalar $[\varepsilon2^r]$, with
$\varepsilon\in\{1,-1\}$ and $\#k=2^{2r}$.
For an odd prime $\ell$, put $f=\operatorname{ord}_\ell(2)$.
The sufficient condition becomes, independently of the
sign $\varepsilon$,
\[
 \begin{cases}
  \gcd(r,f)=1,&f\text{ odd},\\
  \gcd(r,f/2)=1,&f\text{ even}.
 \end{cases}                                      \tag{AC.41.2}
\]
This sharpens AC.38, and also applies to nonscalar
Jacobian Frobenius whenever the displayed deck-group
condition holds over the first field of definition of
the chosen cyclic cover.
\end{claim}
\begin{proof}
The hyperelliptic involution $\iota$ fixes $P$ and acts
as minus one on $J_C$. It therefore preserves the
isomorphism class of every cyclic cover. It has a
unique lift $\tau$ fixing the chosen point above $P$:
two such lifts differ by a deck transformation fixing
that point, hence are equal. The lift inverts the
cyclic deck group, and its square is the identity.
Uniqueness also shows that the lift is defined over
$k_d$ and commutes with its arithmetic Frobenius $F$.
This uses the pointed descent; an arbitrary unpointed
twist need not preserve this commutation relation.

Let
\[
 R=\mathbf F_2[X]/\Phi_n(X),\qquad
 W=H^1_{\mathrm{et}}(D_{\overline k},\mathbf F_2)_{\mathrm{prim}}.
\]
Since $n$ is odd, $R$ is finite etale. Every nontrivial
coherent character space of a genus-two cyclic cover
has dimension one. The geometric etale part therefore
has rank zero or one over every field factor of $R$;
equivalently, $W$ is free of rank one over a direct
factor $R'$ of $R$. This is the same multiplicity
argument used in AC.38. The support $R'$ is stable
under $X\mapsto X^a$ and $X\mapsto X^{-1}$, because
$F$ and $\tau$ act invertibly on $W$.

The action of $F$ is semilinear for the deck
automorphism $X\mapsto X^a$ (or its inverse under the
dual cohomology convention). Since $a^m=1$, $F^m$ is
$R'$-linear, and hence is multiplication by a unit
$b\in R'$. Commuting $F^m$ with $F$ shows that
$b$ is fixed by $a$. Commuting $F^m$ with $\tau$
shows that $b$ is fixed by inversion. Thus
\[
                         b\in(R'^H)^\times.           \tag{AC.41.3}
\]
This argument does not assume that the primitive
factor is ordinary; rank zero factors have simply
been omitted from $R'$.

Over $\overline{\mathbf F}_2$, the primitive character
set is $U$. The field factors of $R$ are the orbits
of multiplication by two, that is, the cosets of $T$.
Taking $H$-invariants groups these into cosets of $TH$.
The stabilizer $T\cap H$ acts on each corresponding
field of degree $|T|$, leaving a field of degree
$|T|/|T\cap H|$. This proves (AC.41.1). A stable
direct factor $R'$ merely removes some of these
field factors.

If $T\subset H$, the fixed ring is a product of copies
of $\mathbf F_2$, whose unit group is trivial. Hence
$b=1$ and $F^m$ is the identity on $W$. The same holds
on the dual geometric two-torsion points of the
primitive abelian factor. A nonzero such point then
remains rational over every further constant-deck
extension. Conversely, geometric two-rank zero
excludes rational two-torsion over every field, so
its class number is always odd. Any field over which
this pointed geometric cover has constant decks
contains $k_{dm}$: first the isomorphism class must
be stable, requiring a multiple of $d$, and then
the deck scalar must become the identity, requiring
a further multiple of $m$.

For scalar Jacobian Frobenius $[s]$ with
$s=\varepsilon2^r$ and $s^2=\#k$, every cyclic
isomorphism class is stable over $k$ and the pointed
deck scalar is $s$. Equivalently, on its Kummer line
the scalar is $s$, and Cartier duality gives the
deck scalar $(\#k)/s=s$. Therefore
$H=\langle2^r,-1\rangle$. For prime $\ell$, if
$f$ is odd then $-1\notin T$, and
$H\cap T=\langle2^r\rangle$; the fixed-field degree
is $\gcd(r,f)$. If $f$ is even, $-1=2^{f/2}$ in
$T$, so $H=\langle2^{\gcd(r,f/2)}\rangle$ and
the fixed-field degree is $\gcd(r,f/2)$. This proves
(AC.41.2).

For composite $n$ the same inclusion and fixed-ring
formulas remain valid. The simple prime formula
must then be replaced by the actual subgroup
calculation in $(\mathbf Z/n)^\times$; in particular,
an even order of two does not by itself imply
$-1\in\langle2\rangle$.
\end{proof}

\paragraph{Scope for the counterexample pencil.}
For a member of the supersingular pencil over
$\mathbf F_{2^{38}}$, the scalar is $-2^{19}$.
If it has no nontrivial cyclic primitive-zero cover
as in AG.33--AG.34, (AC.41.2) excludes every rational
prime-cyclic \textup{(R)} witness whose prime $\ell$
satisfies $19\nmid\operatorname{ord}_\ell(2)$,
over every finite extension. In particular, every
such witness has $\ell\ge191$: for $\ell<191$,
divisibility by19 would force $\ell=19j+1$ with
$j\in\{2,4,6,8\}$, and the four values
$39,77,115,153$ are composite. This is a necessary
condition on witnesses, not a claim that degree191
succeeds or that the arithmetic criterion A3 fails.



\begin{claim}{AC.42}{Proved obstruction to a finite-type-only return bound}
Fixed algebraic coefficient functions evaluated at prime-to-$p$ torsion
points, even when their Frobenius periods are prime to $p$, do not by
themselves give a bounded valuation for the return unit minus one.
More precisely, put $p=2$, let $\omega$ be a primitive cube root of
unity, and set $K=\mathbf Q_2(\omega)$. For $n\ge1$ choose a primitive
$3^n$-th root $\zeta_n$ satisfying
\[
 m_n=3^{n-1},\qquad \zeta_n^{m_n}=\omega.
\]
These roots are prime-to-two Teichmuller points. The fourth-power
Frobenius orbit of $\zeta_n$ has the odd length $m_n$. For the single
fixed degree-one coefficient function $A(X)=1-2X$, its orbit return is
\[
 U_n=\prod_{j=0}^{m_n-1}A(\zeta_n^{4^j})
     =N_{K(\zeta_n)/K}(1-2\zeta_n)
     =1-2^{m_n}\omega.                              \tag{AC.42.1}
\]
In particular
\[
 v_2(A(\zeta_n)-1)=1,\qquad v_2(U_n)=0,
 \qquad v_2(U_n-1)=m_n=3^{n-1}.                    \tag{AC.42.2}
\]
For the full square-Frobenius orbit over $\mathbf Q_2$, the analogous
return is
\[
 N_{\mathbf Q_2(\zeta_n)/\mathbf Q_2}(1-2\zeta_n)
 =\Phi_{3^n}(2)=1+2^{m_n}+2^{2m_n},                 \tag{AC.42.3}
\]
again a unit whose difference from one has valuation $m_n$.
Consequently the annihilator sizes $2^{v_2(U_n-1)}$ grow faster than
$3^{gn}$ for every fixed positive integer $g$.

There is an integral semilinear module realizing these annihilator
sizes. Let $L_n=K(\zeta_n)$, put $M_n=\mathcal O_{L_n}$, and let
$\sigma_n$ be the fourth-power Frobenius of $L_n/K$. Give $M_n$ the
operator
\[
 \mathcal F_n(x)=(1-2\zeta_n)\sigma_n(x).
\]
It normalizes the odd cyclic deck action $\delta_n(x)=\zeta_nx$ by
$\mathcal F_n\delta_n\mathcal F_n^{-1}=\delta_n^4$, and
\[
 \mathcal F_n^{m_n}=U_n,\qquad
 M_n/(1-\mathcal F_n)M_n
       \simeq\mathcal O_K/(2^{m_n}),\qquad
 \exp\bigl((L_n/M_n)^{\mathcal F_n}\bigr)=2^{m_n}.
                                                        \tag{AC.42.4}
\]
This is an algebraic Frobenius-cocycle countermodel to a proposed
intermediate bound. It is not asserted to arise from the Jacobians
in AR.20, or from pure cohomology of Weil weight one. It therefore
does not refute the conditional geometric route in AR.20.
Indeed $|U_n|^2=4^{m_n}+2^{m_n}+1>4^{m_n}$ under a complex
embedding, so the required weight-one equality already fails.
\end{claim}
\begin{proof}
The identity $v_3(4-1)=1$ and the elementary binomial valuation formula
show that $4$ has order $3^{n-1}$ modulo $3^n$. Its powers exhaust
the subgroup of units congruent to one modulo three. Thus the roots
\[
 \{\zeta_n^{4^j}:0\le j<m_n\}
\]
are exactly the roots of $X^{m_n}-\omega$. All prime-to-two roots of
unity belong to unramified extensions of $\mathbf Q_2$, and the
Frobenius over its unramified quadratic extension $K$ is fourth
power. The field degree $[K(\zeta_n):K]$ is therefore $m_n$.
Evaluating the displayed monic polynomial at $X=1/2$ and multiplying
by $2^{m_n}$ proves (AC.42.1). Since $\omega$ and $\zeta_n$ are
two-adic units, (AC.42.2) follows immediately.

The two conjugates over $\mathbf Q_2$ are
$1-2^{m_n}\omega$ and $1-2^{m_n}\omega^2$. Their product is
\[
 (1-2^{m_n}\omega)(1-2^{m_n}\omega^2)
   =1+2^{m_n}+2^{2m_n}.
\]
This is also $\Phi_{3^n}(2)$ by
$\Phi_{3^n}(X)=1+X^{m_n}+X^{2m_n}$. Its difference from one
is $2^{m_n}(1+2^{m_n})$, so the asserted valuation is exact.
Finally $3^{n-1}/n\to\infty$, proving the comparison with $3^{gn}$.

For the integral statement, $M_n$ is free of rank $m_n$ over
$\mathcal O_K$. After extension of scalars to $L_n$, its $m_n$
embeddings diagonalize multiplication, and $\sigma_n$ cyclically
permutes them. Thus $\mathcal F_n$ becomes a weighted cyclic
permutation, with product of its weights $U_n$. Its characteristic
polynomial over $K$ is $T^{m_n}-U_n$, so
\[
                 \det(1-\mathcal F_n)=1-U_n.
\]
Modulo two the operator $\mathcal F_n$ is just the fourth-power
automorphism of $\mathbf F_{4^{m_n}}$. Its fixed space has dimension
one over $\mathbf F_4$. The Smith normal form over the discrete
valuation ring $\mathcal O_K$ therefore has exactly one nonunit
invariant factor, whose valuation equals $v_2(1-U_n)=m_n$.
This proves the cokernel assertion. Applying $1-\mathcal F_n$ to
$0\to M_n\to L_n\to L_n/M_n\to0$, and using its invertibility on
$L_n$, identifies that cokernel with $(L_n/M_n)^{\mathcal F_n}$.
Since $\mathcal O_K$ is unramified quadratic, this finite abelian
group is $(\mathbf Z/2^{m_n}\mathbf Z)^2$.

The coefficient $A$ is invertible at every point under consideration
and has constant reduction one. Thus bounded degree of $A$, bounded
valuation of $A-1$, and periods prime to two all hold in this model.
The growing quantity comes from taking a norm over an increasingly
long orbit. A fixed-locus torsion-distance theorem cannot be applied
to these norm polynomials as though their degrees were bounded.
The exact identities for $1\le n\le6$ are independently checked in
$\mathbf Z[X]/(\Phi_{3^n}(X))$ by
\path{frobenius_norm_obstruction.py}.
\end{proof}


\begin{claim}{AC.43}{Proved limit of a rank-one overconvergence approach}
Let $U/\mathbf F_q$ be smooth and geometrically connected, and let
$\mathcal L$ be a rank-one unit-root overconvergent $F$-isocrystal
on $U$, with coefficients in a fixed finite extension of
$\mathbf Q_p$. Write $b_x$ for its Frobenius eigenvalue at a closed
point $x$. For every fixed integer $c\ge0$, the finite valuations
\[
 \{v_p(b_x-1):x\in|U|,\ v_p(\deg x)\le c,\ b_x\ne1\}
                                                        \tag{AC.43.1}
\]
are uniformly bounded. This assertion would apply to a character
family only after its return eigenvalues had been identified with
such an overconvergent object.

There is a concrete obstruction to that identification in the
scalar example $C_0:y^2+y=x^5$ over $\mathbf F_{16}$. Its ordinary
character unit roots at the order-three parameters and the order-thirteen
parameter of AC.31 cannot be interpolated by one rank-one coefficient
object with finite geometric monodromy on a common geometrically
connected base over $\mathbf F_{16}$ containing those parameter
points with their indicated closed-point degrees two and three.
In particular no rank-one
unit-root overconvergent $F$-isocrystal can have all those prescribed
return eigenvalues. This conclusion concerns the proposed interpolation;
it does not assert that an intrinsic arithmetic $F$-isocrystal on
the ordinary character locus has been constructed.
\end{claim}
\begin{proof}
By \cite[Proposition~8.6, p.~42]{AC_HartlPal2025}, a tensor power
$\mathcal L^{\otimes N}$ is constant. If its constant Frobenius
eigenvalue is $c_0$, then
\[
                         b_x^N=c_0^{\deg x}.           \tag{AC.43.2}
\]
All these quantities are $p$-adic units. Suppose first that $c_0$
is not a root of unity. In a finite splitting field, let $d$ be the
order of its residue. If $d\nmid m$, then
$v_p(c_0^m-1)=0$. Otherwise write $m=dp^a u$ with $p\nmid u$.
The factorization of $X^u-1$ at
$X=c_0^{dp^a}\equiv1$ gives
\[
 v_p(c_0^m-1)=v_p(c_0^{dp^a}-1).
\]
When $v_p(m)\le c$, the latter runs through a finite set of finite
values. Moreover
\[
 v_p(b_x-1)\le v_p(b_x^N-1)=v_p(c_0^{\deg x}-1),
\]
because the other factor of $b_x^N-1$ is integral. If $c_0$ is a
root of unity, (AC.43.2) puts every $b_x$ in one fixed finite set
of roots of unity, which also proves (AC.43.1). This establishes
the asserted conditional bound.

We now prove the obstruction using two genuine character fibers.
On $J_{C_0}$, Frobenius over $\mathbf F_{16}$ is $[-4]$.
Every nonzero point $a_3\in J_{C_0}[3]$ therefore has degree two.
For the Weierstrass point $P=\infty$, the Raynaud theta curve is
$C_0-P$ and contains no nonzero three-torsion: a function with
divisor $3Q-3P$ would have pole order three at $P$, whereas
$L(3P)=\langle1,x\rangle$. The corresponding pointed cubic
cover over $\mathbf F_{256}$ consequently has ordinary new part,
by the rank-one character criterion. The pointed hyperelliptic
involution commutes with Frobenius and interchanges the two
nontrivial cubic characters. Their common polynomial thus has
coefficients in
$\mathbf Z[\zeta_3]\cap\mathbf R=\mathbf Z$. It is
\[
                         1+s_3T+256T^2,
 \qquad s_3\in\mathbf Z,
\]
and has an ordinary unit root $\beta$ satisfying
$\beta^2+s_3\beta+256=0$. The Weil bound and ordinariness imply
$s_3^2<1024$, so $\mathbf Q(\beta)$ is imaginary quadratic.

Let $a_{13}$ be the explicit order-thirteen parameter of AC.31.
Its degree over $\mathbf F_{16}$ is three. A character unit root
$\alpha$ of its pointed cover over $\mathbf F_{4096}$ satisfies
\[
 P(X)=X^4+139X^3+12473X^2+569344X+16777216=0.       \tag{AC.43.3}
\]
More precisely, put $K_0=\mathbf Q(\sqrt{13})$. Its character
quadratic is
\[
 X^2-tX+4096,
 \qquad t=\frac{-139\pm13\sqrt{13}}2,
\]
with discriminant
\[
 d=t^2-16384=\frac{-22009\mp1807\sqrt{13}}2,
 \qquad
 N_{K_0/\mathbf Q}(d)=110486961=39^2\cdot17\cdot4273.
                                                        \tag{AC.43.4}
\]
Both real conjugates of $d$ are negative, so $d$ is not a square
in the real field $K_0$. Also $t=\alpha+4096/\alpha$ belongs to
$\mathbf Q(\alpha)$ and is irrational. Hence
$[\mathbf Q(\alpha):\mathbf Q]=4$ and $K_0$ is a subfield.
The rational number in (AC.43.4) is neither a rational square nor
thirteen times a rational square: its valuation at seventeen is
one. A rational number is a square in $\mathbf Q(\sqrt{13})$
only in those two square classes. Thus $N(d)$ is not a square
in $K_0$. In particular the nontrivial automorphism of $K_0$
cannot extend to $K_0(\sqrt d)$, since such an extension would
make $\overline d/d$, and hence $d\overline d$, a square in $K_0$.
This proves that $\mathbf Q(\alpha)$ is not Galois. It has no
quadratic subfield other than $K_0$: two different quadratic
subfields would generate a biquadratic degree-four field.

Complex conjugation in this CM field fixes $K_0$ and sends
$\alpha$ to $4096/\alpha$. Under the chosen two-adic embedding,
\[
             v_2(\alpha)=0,
 \qquad v_2(4096/\alpha)=12.
\]
No positive power of $\alpha$ can therefore belong to $K_0$.
Since any proper subfield of a degree-four field has degree one
or two, every positive power of $\alpha$ has degree four.

If one rank-one object with finite geometric monodromy had the
specified eigenvalues at these degree-two and degree-three points,
some tensor power would be constant, giving
\[
                     \beta^N=c_0^2,
 \qquad \alpha^N=c_0^3,
 \qquad \alpha^{2N}=\beta^{3N}.
\]
The right-hand side has degree at most two over $\mathbf Q$,
whereas the left-hand side has degree four. This contradiction
proves the interpolation obstruction. For an overconvergent
rank-one unit-root object, the required finite geometric
monodromy follows from the theorem already cited.

The distinction between a full geometric cohomology object and
its unit-root part is essential. The ordinary Legendre family
has overconvergent full cohomology but a unit-root subobject
which is not overconvergent; see
\cite[Theorem~5.4.3.5, Corollary~5.4.3.10 and
Remark~5.4.3.11, pp.~49--52]{AC_Etesse2010}.
In that example the Frobenius coefficient alone can be
overconvergent while its connection is not. Therefore neither
geometric origin nor a finite algebraic formula for the Cartier
coefficient supplies the missing overconvergence hypothesis.
\end{proof}

% AC.44: full Frobenius-packet obstruction for the explicit A6 pencil.
\subsection{Two remaining conductors after simultaneous packet constraints}

\begin{claim}{AC.44}{Proved by an exact finite certificate}
\label{clm:ACpacketconductors}
For any member, with \(t\notin\mathbb F_4\), of the genus-two pencil of
Claim~\ref{clm:AGprimitivecounterexample} over
\(\mathbb F_{2^{38}}\), a nontrivial odd conductor of a
primitive two-rank-zero cyclic cover must belong to
\[
                              \{191,229\}.
\tag{AC.44.1}
\]
Consequently the incidence estimate of Claim~AG.34 improves to
\[
 \#\{\text{bad parameters}\}
       \le133086336+229171540=362257876,
\]
and at least
\[
              274877906944-4-362257876=274515649064
\tag{AC.44.2}
\]
parameters have no nontrivial odd primitive-zero cover.
These are counts of parameters, not of isomorphism classes.
The assertion retains the degree-one qualification of the A6
counterexample and does not exclude rational class-number witnesses.
\end{claim}

\begin{proof}
Let \(n\) be a candidate conductor, and write
\[
 T=\langle2\rangle,\qquad
 U=\langle2,-1\rangle,\qquad
 H=\langle2^{19},-1\rangle
       \subseteq(\mathbb Z/n)^\times.
\tag{AC.44.3}
\]
The theta curve contains the origin, is symmetric, and is stable
under the scalar Frobenius \([-2^{19}]\). Therefore a primitive
theta point \(aQ\), for a chosen generator \(Q\) of the cyclic
subgroup, forces the entire packet \(aHQ\) to lie on theta.
Primitive two-rank zero requires a theta point in every
multiplication-by-two orbit of primitive points. By symmetry it
thus requires at least one \(H\)-packet in each coset of \(U\).
This formulation also handles \(n=191\), where \(-1\notin T\)
and symmetry joins two different doubling orbits.

For any such choice of packets, the set of residues
\[
                  S=\{0\}\cup\bigcup_{aU} a_{aU}H
\tag{AC.44.4}
\]
must satisfy
\[
 \#\{(x,y)\in S^2:x-y=d\}\le2
                 \qquad(0\ne d\in\mathbb Z/n).
\tag{AC.44.5}
\]
Indeed these pairs give distinct points of the intersection of
theta with its translate by \(dQ\). The translation is nontrivial,
the principal polarization has trivial translation stabilizer,
and the intersection number is two. Additional theta points cannot
repair a violation of (AC.44.5), so selecting just one required
packet per \(U\)-coset gives a necessary test.

The single-packet test of Claim~AG.37 leaves only
\(191,229,457,1257\). At \(n=1257=3\cdot419\), the groups
in (AC.44.3) have \(|U|=418\) and \(|H|=22\), while
\(\varphi(n)=836\). Hence two different packets are required.
Each packet has eleven elements in each nonzero residue class
modulo three, since negation exchanges the classes. Thus (AC.44.4)
has occupancies \((1,22,22)\) modulo three and supplies
\[
                   2\cdot22\cdot21=924
\]
ordered pairs whose difference is a nonzero multiple of three.
There are only \(418\) such differences modulo \(1257\).
Equation (AC.44.5) permits at most \(836\) pairs, a contradiction.

At \(n=457\), one has \(|U|=76\), \(|H|=4\), and six
\(U\)-cosets. Choose one theta point as the generator \(Q\),
which normalizes its packet to \(H\) in the identity coset.
Each other \(U\)-coset contains exactly nineteen possible
\(H\)-packets. The accompanying standard-library program
\texttt{primitive\_packet\_sidon.py} enumerates these packets by
group closure and tests all possible extensions of the normalized
packet. It first discards any candidate already violating
(AC.44.5) with \(\{0\}\cup H\); the five remaining candidate
counts are
\[
                            17,19,17,18,18.
\]
It then performs an exhaustive depth-first search. At each step
it adds every ordered difference involving a new packet and
rejects the branch as soon as a multiplicity exceeds two.
The numbers of surviving partial sets containing respectively
one through six packets are exactly
\[
                         1,17,226,879,262,0.
\tag{AC.44.6}
\]
There is therefore no possible full required set for \(457\).
The JSON certificate retains the complete normalization, packet
indices, rejected first-packet witnesses, and search counts;
the source performs the complete search without a heuristic cutoff.
For \(191\) and \(229\) it also retains compatible necessary
sets, making explicit that this test gives no exclusion there.
This proves (AC.44.1).

The incidence estimate is now needed only at these two conductors.
The minimum numbers of primitive theta points are respectively
\(10\) and \(12\), and the numbers of exact-order points in
the ambient product of elliptic curves are
\(J_4(191)=1330863360\) and \(J_4(229)=2750058480\).
Each ambient nonzero torsion point lies on at most one pencil
member. Dividing these two incidence bounds by \(10\) and
\(12\), respectively, gives (AC.44.2). No individually chosen
parameter is asserted to have been certified by this counting
argument.
\end{proof}


\subsection{The fixed norm module and full arithmetic sections}

\begin{claim}{AC.45}{Proved: section origin is weaker than point origin}
\label{clm:ACnormsectionorigin}
Let $C/k$ be smooth, proper and geometrically connected of genus
at least two over a finite field of characteristic $p$, and let
$P\in C(k)$. Use the full etale fundamental groups in the
first row below, and put
\[
 \Omega=\pi_1(C_{\bar k},\bar P),\qquad
 \Delta=\Omega^{(p')},\qquad
 N=\ker(\Omega\longrightarrow\Delta),\qquad
 V=N^{\mathrm{ab},p}.
                                                        \tag{AC.45.1}
\]
Here $N^{\mathrm{ab},p}$ is the maximal abelian pro-$p$
quotient. The actual point section $s_P$ acts on these
groups; write $F$ for arithmetic Frobenius and also for
its induced automorphism of $V$. For the pointed norm
module $M_P$ of AR.11 there is a natural exact sequence
of compact abelian pro-$p$ groups
\[
 0\longrightarrow V\xrightarrow{\ F-1\ }V
      \xrightarrow{\ \varepsilon\ }M_P^F
      \longrightarrow0.
                                                        \tag{AC.45.2}
\]

Every $m\in M_P^F$ is the compatible $p$-adic
abelianized difference from $s_P$ of a continuous section
of the \emph{full} arithmetic etale fundamental group
of $C$. That section projects to exactly $s_P$ in the
geometrically prime-to-$p$ quotient. If $m\ne0$, A1
implies that this full section is not geometrically
conjugate to the section of any point of $C(k)$.

Projection is surjective at every $F$-stable pointed
level, in particular
\[
                    M_P^F\twoheadrightarrow J_C(k)[p^\infty].
                                                        \tag{AC.45.8}
\]
Thus nonzero section-origin classes actually exist
whenever $J_C(k)[p^\infty]\ne0$.

The statement supplies no element with free deck orbit
at every level. In particular it neither proves nor
refutes A10(c) or the nontrivial-cover existence
assertion A3.
\end{claim}

\begin{proof}
\emph{1. The integral norm limit.}
Take the open normal subgroups $H\triangleleft\Delta$
that are stable under $F$. They are cofinal: every
finite geometric cover has a model over a finite
constant extension, hence its defining subgroup has
a finite Frobenius orbit, whose intersection is a
stable normal refinement. Let $U_H\subset\Omega$
be the inverse image of $H$. The point section
gives a model $D_H/k$ with compatible rational lift
$P_H$. Write
\[
                    T_H=T_pJ_{D_H}.
\]
The Abel map based at $P_H$ identifies
$U_H^{\mathrm{ab},p}$ with $T_H$. This comparison
includes the characteristic prime: see
\cite[Proposition~9.1, Lemmas~9.4--9.5 and
Corollary~9.6]{AC_MilneJac1986}.
If $H'\subset H$, the map induced by
$U_{H'}\hookrightarrow U_H$ is the norm map on
Tate modules. Indeed the pointed Abel maps commute
with the morphism of curves and its norm homomorphism
on Jacobians. Equivalently, its dual is pullback on
first etale cohomology, not the transfer multiplied
by the covering degree.

There is an $F$-equivariant isomorphism
\[
                 V\xrightarrow{\sim}\varprojlim_H T_H.
                                                        \tag{AC.45.3}
\]
Here is an elementary continuity verification. A
continuous character $c:N\to\mathbf Q_p/\mathbf Z_p$
has finite image. Choose an open normal subgroup
$O\triangleleft\Omega$ with $N\cap O\subset\ker c$.
Since $[O,N]\subset O\cap N$, the formula
\[
               \widetilde c(no)=c(n),\qquad n\in N, o\in O,
\]
defines a character on the open subgroup $NO$.
Its quotient $\Omega/(NO)$ has prime-to-$p$ order,
since it is a finite quotient of $\Delta$.
Restricting to a sufficiently small $F$-stable normal
subgroup containing $N$ therefore extends $c$ to
one of the $U_H$. Conversely, a character on $U_H$
vanishing on $N$ factors through the pro-prime-to-$p$
group $H$ and is zero. Thus
\[
 \operatorname{Hom}_{\mathrm{cont}}(N,\mathbf Q_p/\mathbf Z_p)
  =\varinjlim_H
    \operatorname{Hom}_{\mathrm{cont}}(U_H,\mathbf Q_p/\mathbf Z_p).
\]
Pontryagin duality gives (AC.45.3), with the transition
maps already identified as norms. No comparison with
the full prime-to-$p$ configuration fiber is involved.

\emph{2. Frobenius coinvariants and rational torsion.}
For every $H$ there is an exact sequence
\[
 0\longrightarrow T_H\xrightarrow{\ F-1\ }T_H
     \xrightarrow{\ \varepsilon_H\ }
       J_{D_H}(k)[p^\infty]\longrightarrow0.
                                                        \tag{AC.45.4}
\]
To check it, put $W_H=T_H\otimes_{\mathbf Z_p}\mathbf Q_p$.
The quotient $W_H/T_H$ is the group of geometric
etale $p$-primary points of $J_{D_H}$. The fixed
subgroup of $T_H$ is zero: otherwise it would give
a nonzero compatible infinitely $p$-divisible
system in the finite group $J_{D_H}(k)$.
Since $T_H$ has finite rank, $F-1$ is therefore
invertible on $W_H$, and
\[
 \varepsilon_H(a)=(F-1)^{-1}a\pmod {T_H}
                                                        \tag{AC.45.5}
\]
gives (AC.45.4). Its image is precisely the
$F$-fixed subgroup of $W_H/T_H$, which is
$J_{D_H}(k)[p^\infty]$. This is the full $p$-adic
Kummer comparison over a finite field; it does not
treat multiplication by $p$ as an etale isogeny.
The maps in (AC.45.4) commute with norms.

For the cofinal stable levels, being fixed by the
action on the entire pointed inverse system means
exactly that each coordinate is rational over $k$.
Consequently
\[
           M_P^F=\varprojlim_H J_{D_H}(k)[p^\infty].
                                                        \tag{AC.45.6}
\]
This also equips $M_P^F$ with its compact pro-$p$
topology, regardless of the unbounded-exponent issue
for the complete module $M_P$.

We verify the inverse-limit exactness explicitly.
For a compatible family $m_H$ on the right of
(AC.45.4), each fiber
$\varepsilon_H^{-1}(m_H)\subset T_H$ is nonempty
and compact. Any finite collection of norm
compatibility equations can be satisfied by choosing
a preimage at one common refinement and taking its
norms to the finitely many levels in the equations.
The remaining coordinates may be chosen arbitrarily
in their nonempty fibers. The finite-intersection
property gives a compatible lift in
$\varprojlim_H T_H$. This proves surjectivity in
(AC.45.2), without a uniform denominator or an
unjustified assertion about restricted transition maps.

If a compatible family $a_H$ maps to zero, write
$a_H=(F-1)b_H$ using (AC.45.4). Each $b_H$ is
unique because $F-1$ is injective on $T_H$.
Commutation with norms then makes the $b_H$
compatible. This proves the kernel equality in
(AC.45.2). Its left injectivity follows at each
level as well. Using (AC.45.3) completes the proof
of that exact sequence.

The asserted surjectivity of projection has an
additional elementary proof. For a transition
$D_{H'}\to D_H$ over $k$, norm composed with
pullback is multiplication by its degree $d$.
Since $d$ is prime to $p$, multiplication by $d$
is an automorphism of $J_{D_H}(k)[p^\infty]$.
Norm on these rational $p$-primary groups is
therefore split surjective. Their finite inverse
system has surjective projection to every
coordinate, proving (AC.45.8). No assertion that
freeness survives these projections is made.

\emph{3. Realization by a full arithmetic section.}
Let $m\in M_P^F$. Choose $v\in V$ with
$\varepsilon(v)=m$, and then choose $n\in N$
mapping to $v$. In the full arithmetic group
\[
        1\longrightarrow\Omega\longrightarrow\pi_1(C,\bar P)
          \longrightarrow G_k\longrightarrow1,
\]
let $F=s_P(\operatorname{Frob}_k)$ and form
$nF$. Every element of a profinite group defines
a continuous procyclic power homomorphism from
$\widehat{\mathbf Z}$. The power homomorphism for
$nF$, followed by the projection to
$G_k\simeq\widehat{\mathbf Z}$, is the identity:
it sends the dense generator $1$ to $1$.
It therefore defines a continuous section $s_n$.
Its prime-to-$p$ geometric projection is exactly
$s_P$, since $n\in N$.

At the pointed level $D_H$, the element $n$ belongs
to $U_H$. With the additive convention for
$T_H$, the abelianized cocycle difference satisfies
\[
      (s_n-s_P)(\operatorname{Frob}_k)=a_H,
                                                        \tag{AC.45.7}
\]
where $a_H$ is the image of $v$ in $T_H$.
Our Kummer convention is
$\delta(x)(F)=F(\widetilde x)-\widetilde x$;
thus its cohomology class is
$a_H\bmod(F-1)T_H$, corresponding to $m_H$
under (AC.45.5). This fixes the sign and verifies
the asserted compatible section realization.

\emph{4. Point origin and the centralizer.}
Suppose $m\ne0$ and $s_n$ were geometrically
conjugate to the section of $Q\in C(k)$.
Projection to the geometrically prime-to-$p$
group and A1 imply $Q=P$.
If $\omega\in\Omega$ conjugates $s_n$ to $s_P$,
its image in $\Delta$ centralizes $s_P(G_k)$,
because these two sections have identical
prime-to-$p$ projections. That geometric
centralizer is trivial by
\cite[Lemma~1.7]{AR_ST2009}. Hence $\omega\in N$.
Its conjugation changes the abelianized difference
by an element of $(F-1)V$, so (AC.45.2) would
give $m=0$, a contradiction.

Thus nonzero norm-module classes have actual full
section realizations that are not point sections.
The argument says nothing further about their
deck stabilizers. In particular it does not replace
the missing point-origin assertion by a freeness
assertion or manufacture a counterexample to A3.
\end{proof}

\paragraph{Relation to the cuspidal proof.}
The point-recognition argument in AG.40 starts from
finite neighborhoods that are known to contain
actual rational points. A section $s_n$ constructed
above does not carry that property merely because
it is a continuous splitting. Consequently its
existence supplies no input to that recognition
step. The exact sequence distinguishes the
Frobenius-fixed Tate module, which is zero, from
the Frobenius coinvariants, which form the
rational torsion norm module. Confusing those two
groups would incorrectly remove the A10 gap.

\paragraph{A small actual-curve example.}
Over $\mathbf F_2$, take the smooth proper model of
\[
                   C:\quad y^2+y=x+\frac1x+\frac1{x+1}.
                                                        \tag{AC.45.9}
\]
The three simple poles $0,1,\infty$ give genus two
and two-rank two by the Artin--Schreier genus and
Deuring--Shafarevich formulas. Each pole supplies
one rational ramification point. There are no
other $\mathbf F_2$-points. Over $\mathbf F_4$,
the two remaining possible $x$-values give
Artin--Schreier right sides of absolute trace one,
so there are still exactly three points. Hence
\[
 \#C(\mathbf F_2)=\#C(\mathbf F_4)=3,
 \qquad P_C(T)=1-T^2+4T^4,
 \qquad \#J_C(\mathbf F_2)=4.
                                                        \tag{AC.45.10}
\]
The exact elementary program
\texttt{norm\_module\_small\_curve.py} verifies
both point counts and the polynomial. In fact
$J_C(\mathbf F_2)\simeq(\mathbf Z/2)^2$:
writing the three ramification points as
$W_0,W_1,W_\infty$, the divisors of $x$ and
$x+1$ are $2(W_0-W_\infty)$ and
$2(W_1-W_\infty)$. The corresponding classes
are distinct and nonzero, because a difference
of distinct points on a positive-genus curve
cannot be principal.

For any choice of rational $P$, the fixed module
$M_P^F$ therefore surjects onto a group of order
four and has nonzero classes represented by
non-point-theoretic full sections. All three
actual point sections are distinct by A1.
No all-level free-orbit property is claimed for
these nonzero module classes.


\paragraph{A fixed geometric curve defeats every finite prime list.}
Let $(C,P)$ be the explicitly specified characteristic-two
curve with its point at infinity in
Claim~\ref{clm:AGexplicitprimitivecounterexample}.
For every finite set $S$ of odd primes there is a finite
extension $L$ of its definition field such that no
$P$-split cyclic cover of degree $\ell\in S$ satisfies
(R) over $L$, or over any further finite extension.
Indeed, there are only finitely many geometric cyclic
covers of each fixed prime degree. Since no nontrivial
odd primitive cover of this curve has primitive two-rank
zero, the new abelian variety of each prime-degree cover
has a nonzero geometric two-torsion point. Choose one
such point for each cover. A common finite extension
makes all the covers and their deck groups defined,
makes a selected point above $P$ rational, and makes
all the selected new two-torsion points rational.

Over this field every geometric cover has its unique
$P$-split twist: constant cyclic twists are parametrized
by the Frobenius translation on the fibre over $P$, and
exactly the zero translation splits that fibre.
The chosen rational new two-torsion persists after every
further extension. Thus every cover in the finite list
has even relative class number throughout those
extensions. Conversely any $P$-split cover of one of
these degrees is the unique pointed form of a member
of the same finite geometric list, so none is omitted.
Consequently there is no finite list of witness primes
that works uniformly over all finite definition fields
even for this single geometric pointed curve. This
statement permits witnesses at other primes and does
not refute A3.
